% MSc dissertation example file, February 2022
%
% Leave one of the documentclass lines uncommented to match your degree.
% You may remove the logo option if it causes problems.
% Do not change any other options.
% \documentclass[logo,msc,adi]{infthesis}     % Adv Design Inf
% \documentclass[logo,msc,ai]{infthesis}      % AI
% \documentclass[logo,msc,cogsci]{infthesis}  % Cognitive Sci
% \documentclass[logo,msc,cs]{infthesis}      % Computer Sci
% \documentclass[logo,msc,cyber]{infthesis}   % Cyber Sec
% \documentclass[logo,msc,datasci]{infthesis} % Data Sci
% \documentclass[logo,msc,di]{infthesis}      % Design Inf
% \documentclass[logo,msc,dsti]{infthesis}    % Data Sci TI
% \documentclass[logo,msc,inf]{infthesis}     % Informatics
\documentclass[logo,msc,ai]{infthesis}           % degree unspecified, do not change except to add your degree
%%%%%%%%%%%%%%%%%%%%%%%%
% Understand any problems and seek approval before assuming it's ok to remove ugcheck.
% \usepackage{msccheck}
\usepackage{Original/msccheck} %LAST SUBMISION CHANGE TO ABOVE AND MOVE FILE BACK OUT OF ORIGINAL

% Include any packages you need below, but don't include any that change the page
% layout or style of the dissertation. By including the ugcheck package above,
% you should catch most accidental changes of page layout though.

\usepackage{microtype} % recommended, but you can remove if it causes problems
\usepackage{xcolor} %remove later this is used for helping edit and draft
\usepackage{subcaption} %subcaptions for plots and figures next to each other. 
\usepackage{amsmath}  %Check we are allowed
\usepackage{amssymb} %Check we are allowed
\usepackage{amsfonts} %Check we are allowed
\usepackage{booktabs} %Check we are allowed
\usepackage{graphicx} %Check we are allowed
\usepackage{tabularx}
\usepackage{float}


%%% Colour commands - delete for final draft 
% - Editing Macros (Hierarchy of Importance) -
% 1. FACT CHECK (Critical): Red commands immediate attention for potential inaccuracies.
\newcommand{\fc}[1]{\textcolor{red}{[FactC: #1]}} 
% 2. CITE (High): Orange highlights missing evidence or attribution.
\newcommand{\ct}[1]{\textcolor{orange}{[CITE: #1]}} 
% 3. APPENDIX (Medium): Blue indicates structural moves (shifting text out of the main body).
\newcommand{\app}[1]{\textcolor{blue}{[APP: #1]}} 
% 4. SMALL REWORD (Low): Teal is for minor stylistic edits, mirroring your friend's setup.
\newcommand{\sr}[1]{\textcolor{teal}{[Sreword: #1]}}
% 5. FIGURE (Low): Purple is for figure-related notes (refer to figure, create figure, etc.).
\newcommand{\fig}[1]{\textcolor{purple}{[FIG: #1]}}
% 6. AI redraft (Low): Ai marker suggestions and corrections
\newcommand{\ai}[1]{\textcolor{brown}{[AI: #1]}}
% 7. LIMITATION (Low): Move to limiations 
\newcommand{\limit}[1]{\textcolor{cyan}{[LIM: #1]}}

%=======================================================================================
\begin{document}
\begin{preliminary}

\title{Attribution of forest top-height growth-curve deviations using repeated LiDAR surveys}
%TODO: DRAFT DIFFERENT VERSIONS OF THE TITLE AND PICK AGAIN AT THE END -- retitled 2026-08-20,
% dropped "Physics-informed prediction" lead (PINN is now supplementary, Q3 only, not the
% dissertation's centre of gravity); attribution now leads, matching the actual Q1->Q2->Q3 story.

\author{Shreeya Gorasia}

\date{\today}
%=======================================================================================

\abstract{TODO: FILL IN}

%TODO: writing: Comeback to this section when the entire report is done 
%=======================================================================================

\maketitle

\newenvironment{ethics}
   {\begin{frontenv}{Research Ethics Approval}{\LARGE}}
   {\end{frontenv}\newpage}

\begin{ethics}
This project was planned in accordance with the Informatics Research
Ethics policy. It did not involve any aspects that required approval
from the Informatics Research Ethics committee.

\standarddeclaration
\end{ethics}

%TODO: writing: Comeback to this section when the entire report is done 
%=======================================================================================

\begin{acknowledgements}
Any acknowledgements go here.
\end{acknowledgements}

%TODO: writing: Comeback to this section when the entire report is done 
%=======================================================================================

\tableofcontents
\end{preliminary}

%TODO: check: Comeback to this when entire report is done and check it and make it fits in one page
%=======================================================================================
\chapter{Introduction} \label{chap:introduction}

\section{Motivation and problem context}
% Forestry problem and wider importance
A Forest Digital Twin is a wider long-term goal for the forestry sector: a dynamic digital representation of forest systems that can support monitoring, forecasting and management planning ~ a shift from static, periodic assessment toward adaptive, data-driven planning. The aim is to combine repeated observations, simulation, and machine-learning(ML) and artificial intelligence (AI) to represent forests as changing systems \cite{baskent2026Review_C1_SS, buonocore2022Proposal_C1_SS}. This is increasingly vital as managers face climate-driven risks—such as extreme storms, destructive winds, and severe summer heatwaves that trigger wildfires and alter growth patterns. At the same time, global forestry is shifting toward continuous cover management, requiring more flexible stand-level models \cite{ForestResearch2009Evidence_C1_SS}. 

Sitka spruce (\textit{Picea sitchensis}) is Britain's primary commercial species, covering approximately 50\% of managed forests \cite{ForestResearch2026Sitka_C1_SS} and contributing roughly \pounds1.1 billion annually to Scotland's economy \cite{ScottishForestry2026Future_C1_SS, mason2011Sitka_C1_SS}. The study focuses on Aberfoyle in Trossachs, Scotland, a Sitka spruce-dominated region with 21 years of repeated LiDAR coverage. Featuring varied topography with exposed, windy ridges to sheltered valleys, providings variation that causes different growth trajectories in different areas of the forest. 

% Height, LiDAR and shared growth relationship
Top height is a key measure of stand development, commonly used in forestry models to estimate downstream metrics like timber volume and site productivity \cite{ForestResearch2026Sitka_C1_SS}, and now can be accurately captured at scale using airborne LiDAR scans \cite{suarezminguez2008Practical_C1_SS}. For most tree species, including Sitka spruce, height development over time is modelled by the Chapman-Richards equation \cite{pommerening2015Methods_C2_CR}, a sigmoid growth curve, with its parameters fixed to represent a stand's maximum height potential and its speed of growth. A fixed parameter curve only captures the average population trend, but local trajectories can deviate from this baseline, due to factors like terrain, climate, soil, wind exposure, and active management all critically constrain or enhance a stand's actual development \cite{socha2023Higher_C1_SS, cameron2015Building_C1_SS}. This creates two related questions: whether local departures from this shared curve are associated with identifiable environmental and management conditions, and whether independent attribution methods agree on which conditions matter.

Whether allowing that relationship to instead vary spatially -- rather than being applied uniformly across an entire region, as current forestry practice including the yield-class system does -- captures real local variation a fixed curve cannot, is a further question this dissertation tests directly (Section~\ref{sec:background-prediction}, Section~\ref{sec:background-spatial-eval}). Flexible ML architectures can also capture such relationships more directly than a fixed curve, without assuming a fixed shape at all; deep neural networks (DNNs) and gradient-boosted decision trees (XGBoost), tested here as a supplementary top-height prediction comparison, learn these relationships from observation.

\paragraph*{Research gap and Contribution}
While repeated airborne laser scanning has been applied to model forest height growth and productivity \cite{janiec2023Development_C1_SS}, limited evidence exists on whether local environmental conditions are associated with departures from Sitka spruce's shared height-growth curve, and whether different attribution methods identify consistent conditions. Many current remote-sensing models do not account for spatial dependence between nearby-identical trees, leading to data leakage and inflated performance estimates \cite{Karasiak2022Spatial_C5_EVAL, ploton2020Spatial_C5_EVAL, lynch2025Digital_C4_NN}. Relying on black-box neural networks without explainable AI limits their use for stakeholders who require interpretable insights to support decision-making.

Therefore the focus of this dissertation targets a small component of the Digital Twin pipeline: using attribution techniques to identify which conditions explain local growth deviations, testing whether allowing that relationship to vary spatially captures what a single global relationship cannot, and testing, as a supplementary comparison, whether any of this signal also improves top-height prediction. Geographically Neural Network-Weighted Regression (GNNWR), itself a neural network, leads the spatial-attribution comparison.

\vspace{5 em}
\vspace{-0.5em}

\section{Objectives and Research Questions}

Answering all three questions below first required cleaning and combining repeated forest
inventory and LiDAR scan records with multi-source geospatial rasters (terrain, wind, climate
and soil) into one unified modelling dataset (Chapter~\ref{ch:data}).

\vspace{-0.5em}
\begin{table}[!htbp]
\small
\setlength{\tabcolsep}{4pt}
\renewcommand{\arraystretch}{0.95}

\noindent\makebox[\linewidth][c]{%
\begin{tabularx}{1.08\linewidth}{
    @{}
    >{\raggedright\arraybackslash}p{0.14\linewidth}
    >{\raggedright\arraybackslash}X
    @{}
}
\noalign{\hrule height 0.9pt}

\textbf{Q}
&
\textbf{Objective, Sub-objectives \& Own Contribution}
\\

\noalign{\hrule height 0.9pt}

\textbf{Q1}:
\newline
Persistent
\newline
Curve-Departure
\newline
Attribution
&
\textbf{(1.1)} Identify environmental and management conditions associated with a
plot's persistent (whole-history) departure from the shared growth curve.
\newline
\textbf{(1.2)} Compare attribution across independent model families (Elastic Net,
XGBoost, a linear mixed-effects model) to test whether they agree on which
conditions matter.
\newline
\textbf{Contribution:} Shows which specific conditions are associated with
persistent departures from the shared growth curve, and where independent model
families agree or disagree.
\\

\hline

\textbf{Q2}:
\newline
Spatial
\newline
Attribution
&
\textbf{(2.1)} Test whether letting the environment-deviation relationship vary
spatially (GNNWR) explains structure the global Q1 models leave unexplained.
\newline
\textbf{(2.2)} Assess the reliability of plot-level differences, and distinguish
real environmental trends from noise and target-construction artefacts.
\newline
\textbf{Contribution:} Tests whether spatial variation captures what global
attribution cannot, and tempers how far any single extreme deviation should be
read as a real biological signal.
\\

\hline

\textbf{Q3}:
\newline
Prediction \&
\newline
Physics
\newline
Guidance
&
\textbf{(3.1)} Compare ML baselines (LR, RF, XGBoost, DNN) and a physics-informed
neural network for top-height prediction under spatial cross-validation.
\newline
\textbf{(3.2)} Test whether the attribution signal identified in Q1 and Q2
translates into a predictive advantage.
\newline
\textbf{Contribution:} Establishes which baseline predicts top height best under
spatial holdout, and states plainly what the physics-informed approach did and did
not achieve for this data.
\\

\noalign{\hrule height 0.9pt}
\end{tabularx}%
}

\label{tab:aims_objectives}
\end{table}
\vspace{-0.8em}

%=======================================================================================
%=======================================================================================
\chapter{Related Work and Background} \label{ch:background} ----- %%%%%% must update, but do after results and evaluation and method has been done baed on the new, attribtuion, ai attirbtuion, prediction quesiton order and story) 
This chapter reviews tree growth modelling, statistical and spatial attribution methods, and machine learning and physics-informed prediction. It concludes by highlighting the research gaps this study addresses. \sr{Reordered 2026-08-20 to lead with attribution, matching the dissertation's actual focus -- full Background rewrite deferred, this is a one-sentence fix only.}

\section{Top-height growth and environmental variation}

The \textbf{Chapman-Richards (CR)} equatioon is widely used in forestry to model height using a flexible sigmoidal curve \cite{richards1959Flexible_C2_CR,pommerening2015Methods_C2_CR,manso2021Dynamic_C2_CR}:
\begin{equation}
H(a) = y_{\max}\left(1 - e^{-ka}\right)^p,
\label{eq:background-cr}
\end{equation}
where $H(a)$ is dominant top height at stand age $a$, $y_{\max}$ is the fitted asymptotic height ceiling (typically 40--50\,m in plantation settings, distinct from the biological maximum usually 80\,m; \cite{puliti2023Tree_C2_CR}), $k$ controls the approach rate to the asymptote, and $p$ governs the inflection point and trajectory shape.

When observational data span a narrow age range, isolating $y_{\max}$ directly is difficult because different combinations of $(y_{\max}, k, p)$ produce overlapping curves over intermediate ages. Observational records also carry measurement noise, repeated-survey correlation, and stand disturbance that functional form alone cannot resolve \cite{rennolls1995Forest_C2_CR}.  British forestry use this relationship to measure site productivity defined as General Yield Class (GYC): the peak average yearly timber volume growth over a stand's life ($\mathrm{m^3\,ha^{-1}\,yr^{-1}}$) \cite{edwards1981Yield_C1_SS,ForestResearch2026Sitka_C1_SS}. For Sitka spruce, empirical yield tables solve directly for GYC by inverting the Chapman--Richards equation:
\begin{equation}
\mathrm{GYC}_{\mathrm{SS}} \approx \alpha \cdot \frac{H(a)}{\left(1 - e^{-k a}\right)^p} - \beta,
\label{eq:gyc-ss}
\end{equation}
where $H(a)$ and $a$ are observed top height and age, while $\alpha$, $\beta$, $k$, and $p$ are calibrated constants for the species (full parameter values given in Appendix~\ref{app:gyc-parameters}). 

%TODO: METNION IN METHOD NOT BACKGROUND: Because GYC is calculated directly from top height via \eqref{eq:gyc-ss}, using inventory-derived productivity ratings as predictors of top height would introduce target leakage.

\begin{figure}[!htbp]
  \centering
  \includegraphics[width=0.65\textwidth]{figures/fig_background/CR-curves-yeildclass.png}
  \caption{\small{CR growth curves Equation~\eqref{eq:background-cr}, across General Yield Classes (GYC 2--24; $\mathrm{m^3\,ha^{-1}\,yr^{-1}}$) for Sitka spruce in Aberfoyle. Higher yield classes scale toward greater asymptotic heights ($y_{\max}$), inverted by Equation~\eqref{eq:gyc-ss} to infer productivity from height and age. Curves reflect regional averages, with individual stands deviating under local site conditions.}}
  \label{fig:chapman-richards-illustrative}
\end{figure}

\paragraph*{Environmental drivers of height growth.} Site-conditioned and mixed-effects have been used to adjust growth-curve parameters using local \emph{physiographic} (terrain and topography) and \emph{bioclimatic} covariates \cite{scolforo2013Dominant_C2_CR,socha2021Regional_C2_CR}. For Sitka spruce across the UK and Ireland, height growth is primarily shaped by elevation, wind exposure, soil moisture, and nutrient availability, for example in northern UK, GYC drops by $3.2\text{--}4.0\,\mathrm{m^3\,ha^{-1}\,yr^{-1}}$ per 100\,m rise in elevation, driven mainly by cooler temperatures and stronger winds \cite{worrell1990Productivity1_C1_SS}. In Ireland studies show soil nutrient levels explain over half the variation in site productivity, with growth slowing at both soil moisture extremes and wind exposure limiting growth primarily on high-elevation sites \cite{farrelly2011Sitka_C1_SS}.

Mapping forest growth across landscapes often pairs inventory records with geospatial data. In Wales, early UK applications linked Forestry Commission sub-compartment records with terrain and soil layers via PCA and regression to model Sitka spruce yield class \cite{bateman1998Using_C1_SS}. Modern approaches combine airborne LiDAR canopy heights with multi-source environmental rasters, demonstrating that macro-climate sets broad height potential while local terrain and soil drive stand-level variation \cite{fricker2019More_C3_LID}. These same data families supplement the Forest Research plot records in this study. However, model reliability hinges on layer resolution; coarse data-sources like SoilGrids show poor ground agreement ($R^2 = 0.19$; requiring careful local validation \cite{miller2024Assessing_C2_CR}).

\section{Data-driven and process-informed prediction} \label{sec:background-prediction}
ML models learn patterns directly from data without assuming a fixed curve shape \cite{lee2018Machine_C3_LID}. Random Forest averages independently fitted trees through bagging and random feature selection, to reduce variance \cite{breiman2001RandomForest_C4_NN}. Gradient Boosted Decision Trees (XGBoost), build trees sequentially to correct the residual errors of the previous ones \cite{chen2016XGBoost_C4_NN}. Both approaches handle non-linear relationships and feature interactions well without requiring manual parameter curves. However, forest remote-sensing studies often test these models on single-date snapshots; predicting repeated height growth requires models to generalize to unseen forest blocks, across stand ages, and over multi-year survey intervals.

\paragraph*{Neural models of forest growth.}
Deep neural networks (DNNs) learn non-linear interactions through layered hierarchical representations, optimising their weights through backpropagation to minimise a specified training loss. But these do not inherently preserve properties such as monotonic growth or bounded asymptotes. Previous comparisons show that unconstrained neural networks can reproduce broad sigmoid growth patterns while failing to preserve these properties across age ranges \cite{ercanli2023Comparison_C4_NN, karatepe2022Total_C4_NN}.

\textbf{PINNs} and hybrid models incorporate scientific relationships into the training objective by combining a data loss with a process loss \cite{karpatne2017Theory_C4_NN,raissi2019Physics_C4_NN, karniadakis2021Physics_C4_NN,willard2022Integrating_C4_NN}:
\begin{equation}
\mathcal{L}_{\text{total}} = \mathcal{L}_{\text{data}} + \lambda\mathcal{L}_{\text{process}},
\label{eq:pinn-loss-general}
\end{equation}
where $\mathcal{L}_{\text{data}}$ measures prediction error, $\mathcal{L}_{\text{process}}$ measures disagreement with the selected process relationship, and $\lambda$ balances the two objectives. Process knowledge can enter neural models in several structurally distinct ways, from bias-correcting a process model's output to embedding a differentiable process model directly alongside parameter sub-networks that let covariates drive its parameters rather than treating them as fixed constants \cite{wesselkamp2024Process_C4_NN}; this project's RQ1 architecture combines physics regularisation with parameter sub-networks of this kind (Section~\ref{sec:rq1-method}). Composite loss functions of this kind can produce imbalanced gradients, causing one loss component to dominate optimisation \cite{wang2021Understanding_C4_NN,rathore2024Challenges_C4_NN}, and there is a risk that a strong penalty based on a misspecified population-level Chapman--Richards curve pulls predictions towards an unsuitable average rather than allowing the network to represent local growth differences.


%------------------------------------------------
\section{Spatial Evaluation and Growth Attribution}
\label{sec:background-spatial-eval}

A common problem when workign with spatial data is that nearby plots share similar soils, climate, and management histories, and thus models can cheat when using random train--test splits as models memorise geographic locations rather than generalisable biological relationships \cite{ploton2020Spatial_C5_EVAL, karasiak2022Spatial_C5_EVAL}. To properly test how well models transfer studies stress the need to test them on separate locations, times, or both To properly test how well models transfer previous work state the importance of using spatially and temporally disjoint holdouts \cite{meyer2018Improving_C5_EVAL}.

Beyond raw prediction, research has evaluated methods for attributing growth differences to specific environmental drivers. Global regularised models such as Elastic Net balance sparsity and collinearity by combining $L_1$ and $L_2$ penalties, yielding interpretable coefficients under assumptions of linearity and spatial stationarity \cite{zou2005Regularization_C4_NN}. Flexible tree ensembles like XGBoost capture complex non-linearities and interactions, with TreeSHAP decomposing predictions into additive feature contributions \cite{chen2016XGBoost_C4_NN,lundberg2017Unified_C7_XAI}. However, SHAP shows how the model makes decisions rather than true biological cause and effect, and correlated variables can distort how important each feature looks. A random intercept in a mixed-effects model, of the kind used for RQ2b (Section~\ref{sec:rq2b-method}), can absorb some of this spatial dependence, but only as a single per-group shift; it cannot let an individual predictor's coefficient itself vary from place to place.

Past studies use spatial regression to test whether environmental associations vary across space. Standard Geographically Weighted Regression (GWR) fits a local linear regression at each location, weighting nearby observations by a kernel function, typically Gaussian or bisquare, of their distance to that location \cite{brunsdon1996Geographically_C6_GWR}. GWR's local coefficients are sensitive to the choice of kernel shape and to the bandwidth, usually selected by cross-validation or AIC minimisation; a single bandwidth applied uniformly can be poorly suited to areas where sample density varies, a recognised limitation of the classical method \cite{brunsdon1996Geographically_C6_GWR}. Two alternatives address this differently without directly solving it: spatial lag and error models add one global spatial-dependence parameter rather than letting every coefficient vary locally, and geographically weighted random forests replace the local linear fit with a local ensemble, gaining flexibility but losing a directly interpretable per-location coefficient for each predictor.

Geographically Neural Network-Weighted Regression (GNNWR) instead replaces GWR's fixed, hand-specified kernel with a neural network that learns the spatial weighting function directly from data: given a target location's distances to the training reference points, the network predicts how much each reference point's coefficients should be scaled, rather than applying a pre-chosen kernel shape and bandwidth \cite{du2020Geographically_C6_GWR}. This keeps GWR's core property, one interpretable linear coefficient per predictor at every location, while removing the requirement to pre-specify how weight should decay with distance, allowing a non-monotonic or anisotropic pattern if the data supports it. It is this combination, spatially varying yet still linearly interpretable coefficients, learned rather than assumed, that motivates using GNNWR for RQ3's attribution over classical GWR, a spatial random intercept of the kind already used for RQ2b, or a geographically weighted random forest (Section~\ref{sec:rq3-method}). Other variants address different limitations of the same family: mixed GWR lets some coefficients stay global while others vary locally, where theory suggests only some predictors should have spatially non-stationary effects; multiscale GWR allows each predictor its own bandwidth rather than one shared bandwidth for all coefficients; and geographically and temporally weighted regression (GTWR), with a neural equivalent (GTNNWR), extends the weighting function to decay with time as well as distance, relevant where repeated observations span enough years for the relationship itself to drift. This project uses none of these extensions: every plot in a given fold is weighted by spatial distance alone, and RQ3 fixes one coefficient per predictor per location rather than allowing some to stay global, the simplest member of the family consistent with what Section~\ref{sec:rq3-method} motivates.

In China, a GNNWR was used to model fallen trees from terrain and stand variables, outperforming global models and classical GWR while reducing residual spatial autocorrelation \cite{sun2022Local_C6_GWR}. However, 355 of the 443 plots were randomly selected for training and the remaining 88 for validation within the same reserve. The study therefore showed that learned spatial weights captured variation within the sampled reserve, but they did not test transfer to spatially disjoint areas within that reserve or to a separate forest.

%------------------------------------------------
\section{Research gap and study position}
\label{sec:background-gap}

Established top-height models provide simple representations of stand development, yet uniform height-age curves often mask significant local growth variation. Multi-temporal LiDAR data provides extensive spatial coverage of dominant canopy height, but these metrics reflect a combination of biological growth, sensor acquisition differences, and unmeasured plot-level management.

ML has already been used to relate LiDAR-derived canopy height to environmental conditions in other forests \cite{fricker2019More_C3_LID}. However, no study identified in this review combines repeated LiDAR observations with stand records and GIS-derived terrain, wind, climate and soil variables to model Sitka spruce height trajectories at Aberfoyle, where direct field measurements are limited. ML and PINNs provide flexible alternatives to static parametric curves, but adding domain guidance does not inherently guarantee better predictions than an unconstrained DNN or a tuned tree ensemble. Separately, the attribution of plot-level growth-curve departures has not been compared across linear, tree-based and spatially varying models under spatially independent validation, leaving it uncertain whether different model families identify consistent environmental influences.

This dissertation addresses these gaps through the aim and objectives set out in the introduction.


% Explain what previous studies did.
% State the important findings.
% Compare the approaches.
% Identify limitations that lead directly to the dissertation’s research gap.


%=======================================================================================
%=======================================================================================
\chapter{Data} \label{ch:data}
This chapter describes the data used for model training and evaluation, including cohort construction, feature engineering and the main data limitations. Table~\ref{tab:data_summary} summarises the predictor and target variables, their sources and temporal coverage. The rationale for feature selection, data splitting and model evaluation is given in Chapter~\ref{ch:methodology}.

\section{Study area and cohort coverage}
The dataset covers approximately $16\text{~km} \times 12\text{~km}$ of Aberfoyle and contains 795,645 raw plot observations across six survey years (2002--2023). The records contain airborne laser scanning (ALS)-derived variables, including canopy cover and height percentiles, alongside stand and management fields such as planting year and known felling history (Table~\ref{tab:data-fields}). 
Figure~\ref{fig:data-study-area} shows the study-area topography, the spatial coverage of both cohorts, and an example of standard and boundary-clipped plots.

\begin{figure}[!htbp]
\includegraphics[
    width=1.05\textwidth,
    height=0.48\textheight,
    keepaspectratio
]{figures/fig_data/aberfoyle_area.png}
\caption{\small{Left:Aberfoyle study area. Middle: cohort coverage, Right: plot-level detail}}
\label{fig:data-study-area}
\end{figure}

Due to a LiDAR coverage increase of $37\text{~km}^2$ between 2006 and 2008, two nested cohorts were created (Figure~\ref{fig:cohort_coverage_funnel}): a \emph{4-survey} ($\in \{2008, 2012, 2021, 2023\}$) and a \emph{6-survey} ($\in \{2002, 2006, 2008, 2012, 2021, 2023\}$). The former has broad spatial coverage, while the latter extends temporal trajectories. Sequential filters required complete survey presence, pure Sitka spruce, planting between 1850 and the survey year, ages $0\text{--}200$ years, heights $0\text{--}60\text{~m}$ (ecological max; \cite{mason2011Sitka_C1_SS}), and age $\ge 30$ in 2023 to avoid unreliable young-stand heights (Table~\ref{tab:data-cleaning-funnel} \cite{suarezminguez2008Practical_C1_SS}).

\begin{table}[!htbp]
\centering
\scriptsize
\begin{tabular}{lrrrr}
\toprule
& \multicolumn{2}{c}{4-survey} & \multicolumn{2}{c}{6-survey} \\
\cmidrule(lr){2-3}\cmidrule(lr){4-5}
Stage & Plots & Removed & Plots & Removed \\
\midrule
0. Raw dataset
& 230,359 & -- & 230,359 & -- \\
1. Cohort years present
& 230,345 & 0.01\% & 230,359 & 0.00\% \\
2. Complete survey coverage
& 112,782 & 51.04\% & 18,781 & 91.85\% \\
3. Sitka spruce only
& 72,395 & 35.81\% & 13,897 & 26.01\% \\
4. Valid planting year
& 71,766 & 0.87\% & 13,897 & 0.00\% \\
5. Plausible age (0--200 years)
& 71,766 & 0.00\% & 13,897 & 0.00\% \\
6. Plausible height (0--60 m)
& 71,766 & 0.00\% & 13,897 & 0.00\% \\
7. Mature stands (age $\geq 30$ in 2023)
& \textbf{58,112} & 19.03\%
& \textbf{13,769} & 0.92\% \\
\bottomrule
\end{tabular}
\caption{\small{Plots left after each filter. Percentages show the share removed from the previous row, rather than from the raw dataset. Figure~\ref{fig:data-study-area} shows the final spatial coverage.}}
\label{tab:data-cleaning-funnel}
\end{table}

%------------------------------------------------
\section{Response variable and derived attribution targets}
\label{sec:data-response}
\sr{Retitled from "Target variable and original stand fields" -- the old title and opening
implied one target (prediction). There are three, one per Q, all built from the same base
field.}

The base measurement every target in this dissertation is built from is the plot-level 95th
percentile ALS return height above ground (\texttt{elev\_\allowbreak percentile\_\allowbreak
95th}) per survey. Three distinct targets are derived from it: the raw height itself (Q3,
prediction); the mean deviation from the shared Chapman--Richards growth curve, averaged across
a plot's full survey history (Q1, persistent attribution); and the deviation between a plot's
own fitted asymptote and its yield-class-implied asymptote (Q2, spatial attribution). How the
shared curve is fitted and each derived target is constructed is defined in
Chapter~\ref{ch:methodology}. Records also supply planting year, management history, and Forest
Research yield and volume estimates (Table~\ref{tab:data-fields}). To prevent data leakage, any
fields derived directly from the height target were excluded from prediction inputs.

\begin{table}[!htbp]
\scriptsize

\setlength{\tabcolsep}{3pt}
\renewcommand{\arraystretch}{0.95}
\begin{tabular}{
    >{\raggedright\arraybackslash}p{2.6cm}
    >{\raggedright\arraybackslash}p{7.0cm}
    >{\raggedright\arraybackslash}p{3.2cm}
}
\hline
Field & Meaning or derivation & Role in this study \\
\hline
\texttt{elev\_\allowbreak percentile\_\allowbreak 95th} & 95th percentile of ALS return heights above ground, used without correction & Response variable \\
\texttt{Top\_Height95} & $=\texttt{elev\_percentile\_95th}\times1.1$; Forest Research's own corrected height estimate & Retained for reporting but not modelled \\
\texttt{plyr}, \texttt{Age} & Planting year; \texttt{Age} = survey year $-$ \texttt{plyr} & Input feature; used for data cleaning \\
\texttt{yldc} & An existing Forest Research yield-class field, not calculated from this project's LiDAR-derived height response & Excluded from prediction inputs; used to construct Q2's own attribution target (\sr{not just "a comparison variable" -- it is load-bearing for Q2, reworded to say so}) \\
\texttt{whcl} & Ordinal Forest Research windthrow hazard class, 0-6 & Management descriptor \\
\texttt{CanopyCover} & Canopy-cover fraction from the same ALS survey and processing step as the height response & Input feature \\
\texttt{Thin}, \texttt{last\_thinn}, \texttt{next\_thin\_} & Thinning history; \texttt{Thin} is binary (1 if the stand has been thinned, 0 if never thinned); \texttt{last\_thinn}/\texttt{next\_thin\_} record the last and next thinning years & Input feature \\
\texttt{Vol95}, \texttt{area}, \texttt{g1}, \texttt{g3} & Forest Research volume estimate, calculated from height, plot area and canopy cover & Used for disturbance screening \\
\texttt{GYCspec95}, \texttt{p1-p5} & General Yield Class, calculated directly from a plot's own measured height (\texttt{Top\_Height95}) & \texttt{GYCspec95}: excluded from prediction, calculation incorporates the height response, target leakage. \texttt{p1-p5}: excluded from prediction inputs for the same reason, but \texttt{p4}/\texttt{p5} are deliberately used, held fixed, to construct Q2's yield-class-implied asymptote -- excluded and load-bearing are not the same thing \sr{(this dual role was previously collapsed into one blanket "excluded" -- clarified)} \\
\hline
\end{tabular}
\caption{\small{Target, stand and management fields provided by Forest Research.}}
\label{tab:data-fields}
\end{table}
%------------------------------------------------
\section{Derived terrain, wind, climate and soil variables}
\label{sec:data-environmental}

Each plot was matched with external terrain, wind, climate, and soil variables (Table~\ref{tab:data-env-sources}). Appendix~\ref{tab:data-env-variables_detail_app} defines each variable.

\begin{table}[!htbp]
\centering
\scriptsize
\setlength{\tabcolsep}{4pt}
\renewcommand{\arraystretch}{1.05}

\begin{tabular}{
    >{\raggedright\arraybackslash}p{2.0cm}
    >{\raggedright\arraybackslash}p{4.2cm}
    >{\raggedright\arraybackslash}p{2.3cm}
    >{\raggedright\arraybackslash}p{4.0cm}
}
\hline
Data family & Source and resolution & Time basis & Variables used \\
\hline

Terrain
& OS Terrain 50, 50\,m \cite{ordnancesurveyTerrain50_C3_LID}
& Static
& Elevation, slope, aspect, northness and eastness \\

Soil
& CEH Natural Capital, 50\,m \cite{feeney2026CEHNaturalCapital_C3_LID}; SoilGrids v2.0, 250\,m \cite{poggio2021SoilGrids_C3_LID}
& Static
& Pedotope, drainage, texture and soil pH \\

Wind exposure
& Global Wind Atlas, 250\,m \cite{davis2023GlobalWindAtlas_C3_LID}
& 2006--2015 climate
& Long-term wind speed and exposure \\

Observed wind
& MIDAS Open weather stations \cite{metoffice2019MIDASOpen_C3_LID}
& Survey intervals
& Mean and extreme wind conditions \\

Long-term climate
& CHELSA v2.1, 1\,km \cite{karger2017CHELSA_C3_LID}
& 1981--2010 climate
& Growing degree days and precipitation \\

Matched climate
& HadUK-Grid, 1\,km \cite{hollis2019HadUKGrid_C3_LID}
& Cohort years
& Temperature and frost \\

Spatial context
& OS Open Roads and Open Rivers, vector data
\cite{ordnancesurveyOpenRoads_C3_LID,ordnancesurveyOpenRivers_C3_LID}
& Static
& Distance to roads, rivers and forest edges \\

\hline
\end{tabular}

\caption{\small{External data sources used in the analysis. Values were matched to plot centres. Full source, resolution, coverage and extraction details are given in Appendix~\ref{app:data-sources}.}}
\label{tab:data-env-sources}
\end{table}

The Global Wind Atlas gives long-term modelled wind, TOPEX measures terrain shelter, and MIDAS records wind at three weather stations. CHELSA describes long-term climate from 1981 to 2010 \cite{karger2017CHELSA_C3_LID}. HadUK-Grid temperature and frost values were averaged across each cohort's survey years \cite{hollis2019HadUKGrid_C3_LID}.
% NOTE (2026-08-20, verified against code -- notebooks/environmental_data/av1_aux_data_resolution_check.ipynb,
% cell 50): label is accurate, not misleading. This is genuinely CHELSA's global climate-normal
% product, computed once over 1981-2010 and never updated per-year, downloaded directly from the
% source with the period stamped into the file's own name
% (CHELSA_bio1_1981-2010_V.2.1.tif, from os.zhdk.cloud.switch.ch/chelsav2/GLOBAL/climatologies/1981-2010/bio/).
% Not "data from 1981" in the sense of a single year's observation -- a 30-year statistical
% average (a "climate normal"), which is exactly what makes it a real temporal mismatch against
% the 2008-2023 survey period, not a naming confusion. Fold into prose once placement/phrasing
% decided.


%=======================================================================================
%=======================================================================================
\chapter{Methodology and Experimental Design} \label{ch:methodology}
This chapter describes how the raw data was used to answer three questions, in the
order the Results chapter answers them. It defines the splits, feature sets, models,
targets and evaluation procedures. The same held-out observations were used whenever
models were compared.

\paragraph*{Experimental framework:} Q1 attributes a plot's persistent departure from
the shared growth curve to environmental and management conditions. Q2 tests whether
letting that relationship vary spatially explains what Q1's global models cannot. Q3,
supplementary, tests whether any of this signal also improves top-height prediction,
and reports what a physics-informed architecture did and did not achieve.
\begin{table}[!ht]
\centering
\footnotesize
\renewcommand{\arraystretch}{1.15}

\begin{tabular}{
@{}
>{\raggedright\arraybackslash}p{0.9cm}
>{\raggedright\arraybackslash}p{3.5cm}
>{\raggedright\arraybackslash}p{3.2cm}
>{\raggedright\arraybackslash}p{2.8cm}
@{}
}
\toprule
Q & Unit and target & Models or comparison & Main evaluation \\
\midrule
Q1 & Plot; mean pooled-CR residual & Elastic Net, XGBoost and two-stage
mixed-effects model & Five-fold spatial block CV; coefficients, SHAP and residual Moran's $I$ \\
Q2 & Plot; difference between fitted and yield-class-implied asymptotes & Elastic Net, XGBoost
and GNNWR & Five-fold spatial block CV; attribution, Moran's $I$ and target audit \\
Q3 & Plot--survey row; 95th-percentile top height, and row-level error reduction against the
fold-matched CR curve & CR, linear regression, random forest, XGBoost, DNN, PINN and PINN-$k$
& Five-fold spatial block CV; random-split and PINN forward-pass diagnostics; CR-error-quartile
correction with permutation control \\
\bottomrule
\end{tabular}
\caption{Units, targets and evaluation used for each question. Q3's row-level CR-error-quartile
check reuses Q3's own model predictions (conditioned DNN, PINN, PINN-$k$ and XGBoost) rather
than introducing a separate comparison.}
\label{tab:rq-framework}
\end{table}

\begin{figure}[!htbp]
\centering
\includegraphics[width=\textwidth]{figures/fig_methodology/method_fig_pipeline-overview.png}
\caption{\small{Analysis pipeline. The two cohorts feed shared spatial partitions and separate feature sets. Q1 examines departure from a shared curve, Q2 examines plot-specific asymptote deviation, and Q3 predicts top height.}}
\label{fig:pipeline-overview}
\end{figure}

\section{Data partitions and leakage control}
The four-survey cohort was the main dataset; the smaller six-survey cohort provided a
nested sensitivity analysis with longer observation histories. The unit of analysis
depended on the question. Q1 and Q2 aggregated the data to one observation per plot.
Q3 used plot--survey observations.

\subsection*{Evaluation splits}
Spatial-block cross-validation was the primary evaluation design. Whole forest compartments were assigned to folds, preventing the same compartment from appearing in both training and testing \cite{roberts2017Cross_C5_EVAL}. Training plots within 60\,m of a validation or test plot were also removed, while all held-out plots were retained. Five folds were balanced approximately by row count. In each rotation, one fold was used for testing, the following fold for validation and the remaining three for training. Each compartment was therefore tested once and predictions from the five test folds were pooled before calculating the overall metrics \cite{ploton2020Spatial_C5_EVAL}.

Early exploratory experiments used a random plot-level split (60/20/20 ratio, seed~42) as an easier reference without spatial separation. This split showed clear leakage -- near-duplicate neighbouring plots landing in both training and test -- and substantially inflated apparent accuracy, motivating the spatial-block design used as the primary (and, in this chapter, only reported) evaluation design instead; plot-level results are not carried into the Results chapter. \limit{A temporal split (holding out later survey years) was also constructed for Q3 but is not reported: the chapter's scope was narrowed to the spatial-block comparison and the Q1/Q2 attribution questions, and the temporal holdout was not carried through to a full, comparably-robust analysis. Figure~\ref{fig:splits-combined} should be re-cropped to show only panels A and B before submission.}

\begin{figure}[!htbp]
\includegraphics[width=\textwidth]{figures/fig_methodology/method_fig-splits-combined_draft.png}
\caption{Evaluation data-splitting strategies: (A) random plot-level split and (B) compartment-level spatial-block split with a 60\,m buffer.}
\label{fig:splits-combined}
\end{figure}

Continuous predictors were standardised using parameters estimated from the training partition. The same transformation was applied to validation and test data. Age was scaled separately so that the PINN derivative could be converted to metres per year. Categorical soil classes were one-hot encoded. Compartment, sub-compartment and block identifiers were excluded from the model inputs. Two neighbour-height variables were also excluded because they had been calculated from nearby plots' observed heights before partitioning and could therefore leak held-out outcomes.

%------------------------------------------------
\section{Environmental feature sets}

Because Q1, Q2, and Q3 target different response variables, feature sets are constructed separately for each question using a shared four-step protocol:

\begin{enumerate}
\setlength\itemsep{0pt}
\item \textbf{Coverage screening:} Variables with insufficient spatial coverage are removed prior to ranking.
\item \textbf{Collinearity filtering:} For collinear pairs with Spearman $|\rho| \ge 0.95$, one representative is retained, prioritising direct external measurements over derived indices.
\item \textbf{Composite ranking:} Remaining candidates are sorted by their average rank across univariate Spearman $|\rho|$, permutation feature importance, and drop-column loss in held-out $R^2$.
\item \textbf{Multicollinearity pruning:} Features are pruned iteratively until all variance inflation factors (VIF) fall below 5 \cite{obrien2007Caution_C5_EVAL}. Protected stand variables remain fixed in the matrix throughout this process to constrain the VIF of competing predictors.
\end{enumerate}

Four nested feature sets are defined from these steps. Set~1 contains only the protected baseline variables: canopy cover, time since thinning, its missingness indicator, and recent thinning. Set~2 adds the ten highest-ranked candidates. Set~3 expands on Set~2 by adding the top half of terrain and wind features, while Set~4 adds the top half of climate, soil, and spatial context features to Set~3. Appendix~\ref{app:feature-sets} lists complete variable assignments and VIF diagnostics.

Feature selection is performed once on the designated training partition rather than re-run across individual folds, which ensures all models train on identical feature subsets, meaning fold-to-fold variance reflects true model stability and spatial partitioning rather than shifting predictor selection.

%------------------------------------------------
\section{Q1: Persistent curve-departure attribution}
\label{sec:rq2b-method}
\label{sec:q1-method}

Q1 focuses on the four-survey cohort, as the six-survey cohort contains too few compartments for compartment-level spatial attribution. For each plot-survey observation, the deviation from expected height is:
\begin{equation}
r_{it}=H_{it}-H(a_{it}),
\end{equation}
yielding the time-averaged plot target:
\begin{equation}
\bar r_i=\frac{1}{T_i}\sum_{t=1}^{T_i}r_{it}.
\label{eq:mean-cr-residual}
\end{equation}
Positive values identify stands that consistently outperform the pooled Chapman--Richards curve (the pooled fit itself is defined in Section~\ref{sec:rq1-method}). Stand age is excluded from the candidate predictor matrix because it already defines the baseline expectation.

Elastic Net and XGBoost are evaluated under five-fold compartment cross-validation. Elastic Net tunes its penalty parameter and $\ell_1$ mixing ratio through nested cross-validation on each training partition. XGBoost uses the same fixed tuned configuration as Q2 (tree depth 4, learning rate 0.04, up to 500 boosting rounds, early stopping against a held-out validation split) -- adopted after a sensitivity check showed the raw library defaults it originally ran on gave a real, substantial $R^2$ cost (0.07--0.09 pooled per set); Q3 uses a separately per-target tuned configuration instead (Section~\ref{sec:rq1-method}). Model performance is measured using pooled out-of-fold $R^2$.

To account for hierarchical spatial grouping, a two-stage mixed-effects model is also estimated. Stage~1 computes $\bar r_i$ from the non-linear baseline curve. Stage~2 fits:
\begin{equation}
\bar r_i=\boldsymbol{\beta}^{\top}\mathbf{x}_i+u_{c[i]}+\varepsilon_i,
\label{eq:mixed-effects}
\end{equation}
where $u_{c[i]}\sim\mathcal{N}(0,\sigma_u^2)$ represents a random intercept for each compartment, without random slopes. Fixed-effect coefficients are estimated using restricted maximum likelihood (REML) within each training fold. Spatial variance explained is calculated as the proportional reduction in compartment random-intercept variance relative to an intercept-only baseline; this comparison switches to maximum likelihood (ML) rather than REML for both models, since REML variance-component estimates are not cleanly comparable across models with different fixed-effect sets -- REML profiles out the fixed effects in a way that depends on how many there are, so a naive REML-vs-REML comparison would confound "the fixed effects explained variance" with "the REML transform changed." Parameters are reported as the mean $\pm$ sample standard deviation across folds.

Four diagnostic checks aid interpretation. First, Set~4 is re-estimated without canopy cover. Second, Pearson correlations are calculated between canopy cover and each target variable. Third, SHAP values are recomputed without canopy cover to determine whether correlated predictors absorb its importance. Fourth, between-compartment variance shares are computed for slope and TOPEX to resolve differences between the mixed model and Elastic Net.

Together, Elastic Net coefficients, mixed-effects fixed effects, and XGBoost mean absolute SHAP values provide three distinct views of growth attribution. SHAP calculations are performed post-estimation and do not influence feature selection. Finally, Moran's $I$ is computed on the pooled held-out residuals from Elastic Net and XGBoost to check for remaining spatial autocorrelation (method in Section~\ref{sec:eval-robustness}).

\ai{Proposed additional evaluation, not yet run: every Q1 number above (coefficients, $R^2$, Moran's $I$) reflects a single spatial-block split-seed (42). Q2's GNNWR is already reseeded three times to test sign stability (Section~\ref{sec:q2-method}); repeating that reseed for Q1's Elastic Net and XGBoost, reusing the same five-fold compartment assignments under different seeds, would test whether the slope/TOPEX stability contrast reported in Section~\ref{sec:results-q1-attribution} holds up outside one particular fold assignment, directly answering a gap already flagged in this project's own audit notes rather than leaving it a single-seed number reported as if stable.}

%------------------------------------------------
\section{Q2: Spatial attribution}
\label{sec:rq3-method}
\label{sec:q2-method}

Q1's global attribution models leave real, significant spatial autocorrelation in their residuals across every set, though Elastic Net's own residual clustering barely responds to added environmental data while XGBoost's responds substantially more (Section~\ref{sec:results-q1-attribution}), and a compartment-level random intercept only partially reduces it. This motivates testing whether letting the environment--deviation relationship itself vary spatially, rather than fitting one global relationship and shifting only its intercept locally, captures structure the global models cannot. GNNWR is used for this rather than classical GWR because it removes the need to pre-specify a distance kernel and bandwidth, and rather than a geographically weighted random forest because it keeps one interpretable linear coefficient per predictor at each location, preserving direct comparability with the Elastic Net and mixed-effects coefficients reported for Q1 (Section~\ref{sec:background-spatial-eval}). Keeping the functional form linear is deliberate: it isolates spatial non-stationarity as one clean, testable factor, holding it apart from functional-form flexibility, which Q1 already tested separately via XGBoost.

\limit{GNNWR's coefficients are local but still linear -- one straight-line effect per predictor per location, the same functional-form assumption as Elastic Net, not XGBoost. Q1 and Q3 both find evidence of real non-linear and interaction structure in this data (XGBoost erasing more residual spatial clustering than Elastic Net; a substantial share of XGBoost's SHAP weight sitting in interaction terms), so a location where GNNWR's coefficient looks unstable or small cannot be cleanly attributed to a genuine spatially-varying effect versus a straight line fitting a curved relationship differently well from place to place -- the two look identical from GNNWR's output alone. This does not undermine the R2-based comparison reported for Q2 (functional form is held fixed there specifically to isolate the spatial question), but it does limit how confidently any single local coefficient, taken on its own, can be read as evidence of a real spatially-varying environmental effect.}

Q2 estimates a distinct asymptotic height ceiling for each plot while fixing the Forestry Commission rate and shape parameters. For plot $i$ at survey time $t$, let
\begin{equation}
g_{it}=\left(1-e^{-p_{4,i}a_{it}}\right)^{p_{5,i}}.
\end{equation}
Fitting the single-parameter growth curve through the origin by ordinary least squares yields
\begin{equation}
\hat y_{\max,i}=\frac{\sum_t H_{it}g_{it}}{\sum_t g_{it}^2}.
\label{eq:plot-ymax-fit}
\end{equation}
The expected asymptote implied by standard yield class is calculated directly from the regional Forestry Commission parameters:
\begin{equation}
y_{\max,i}^{\mathrm{yldc}}=\mathrm{yldc}_i p_{2,i}+p_{1,i}+2p_{3,i}.
\end{equation}
The final attribution target is the deviation between these two ceilings:
\begin{equation}
\Delta y_{\max,i}=\hat y_{\max,i}-y_{\max,i}^{\mathrm{yldc}}.
\label{eq:local-ymax-diff}
\end{equation}
This formulation keeps the target strictly free of environmental covariates. Fixing $p_4$ and $p_5$ avoids over-parameterising curves from only four to six time points, though it absorbs any local shape mismatch into the fitted asymptote -- the mechanism behind the extreme-asymptote pattern audited below. Note $p$ here plays a different role from the pooled Chapman--Richards curve's own shape parameter fitted in Section~\ref{sec:rq1-method}: this section's $p_4$/$p_5$ are fixed, externally-sourced Forestry Commission values, not re-fitted quantities, despite the shared notation.

Elastic Net and XGBoost models are trained on $\Delta y_{\max}$ using feature Sets~2--4, with the same training setup as Q1 (Section~\ref{sec:q1-method}). To model spatially varying associations, Geographically Neural Network-Weighted Regression (GNNWR) is also estimated \cite{du2020Geographically_C6_GWR}:
\begin{equation}
\hat y_i=\beta_0(s_i)+\sum_j\beta_j(s_i)x_{ij}.
\label{eq:gnnwr}
\end{equation}
For each target location $s_i$, the spatial weighting sub-network receives a MinMax-scaled vector of Euclidean distances to all training reference plots. The network predicts multiplicative scaling factors that adjust a frozen, global ordinary least squares baseline for the intercept and each slope parameter. All reported runs use the full, uncapped training reference set without predefined distance kernels or fixed bandwidths. Predictions for held-out validation plots reference these same training points, evaluated across the shared five spatial folds.

\subsection*{Implementation constraints and simplifications}
Two separate hardware constraints shaped how GNNWR was fitted, disclosed here because they affect what the reported diagnostics can and cannot show. First, the spatial-weighting sub-network's input width equals the reference-set size, since it takes each plot's full distance vector to every reference point; at this project's scale ($\sim$31,000 training plots) this exhausted the available cluster GPU's memory on the first optimiser step. An initial workaround capped the reference set to a fixed row count, sampled proportionally across compartments so a smaller reference set still covered the forest geographically rather than a lucky or unlucky random slice; this was later superseded by removing a second, unrelated bottleneck (below), which allowed the reference set to be restored to its full, uncapped size for every headline result reported in this chapter. A separate reference-density ablation (Section~\ref{sec:results-q2-spatial}) deliberately re-introduces this cap at several sizes to test whether reference sparsity explains the six-survey cohort's instability; those runs are the one exception to "full, uncapped" and are reported as such. Second, and independently, the underlying GNNWR package's diagnostics class builds a classical-GWR hat matrix to report training-time AIC and F-test statistics, an $O(n^2 \times \text{features})$ tensor that exceeded available GPU memory regardless of reference-set size. This hat-matrix construction was disabled above a row-count threshold. It does not affect the $R^2$/RMSE values reported throughout this chapter, which are computed by a separate, direct formula; it does mean AIC/AICc are approximated by plain feature count rather than the true GWR-corrected effective degrees of freedom, and that the F1/F2/F3 significance tests classical GWR software reports are unavailable here. Neither constraint affects EN or XGBoost, which do not share GNNWR's architecture.

Predictor importance is assessed using both tree split gain and post-hoc SHAP values. XGBoost gain sums the training loss reduction across all tree splits involving a feature, reflecting its utility during model construction. In contrast, mean absolute SHAP values quantify each feature's average additive contribution to final predictions \cite{lundberg2017unified}. SHAP values are computed out of fold across each test partition and pooled before interpretation. A separate check groups predictors into terrain, wind, soil, climate, and spatial-position/edge-effect categories and measures each category's contribution by grouped permutation importance (mean $R^2$ drop when every column in that category is shuffled together) and by the change in residual Moran's $I$ when that category is removed entirely and the model refitted, on a single spatial-block split of Set4 (4survey only).

\subsection*{Target and outlier audit}
To prevent flawed time series from corrupting the analysis, a plot is excluded from the curve-fitting population if any consecutive survey interval meets either of two filtering rules:
\begin{itemize}
\setlength\itemsep{2pt}
\item \textbf{Clearfell pattern:} Top height falls by at least 25\%, canopy cover drops by at least 70\%, and standing volume drops by at least 70\%.
\item \textbf{Measurement inconsistency:} Top height falls by at least 25\%, while canopy cover drops by no more than 10\% and standing volume does not decrease.
\end{itemize}
Intervals exhibiting large height drops that match neither criterion are classified as ambiguous and kept in the dataset.

Unusually low or high fitted asymptotes are screened separately using Tukey's $1.5\,\mathrm{IQR}$ fences \cite{tukey1977Exploratory}:
\begin{equation}
\left[Q_1 - 1.5\,\mathrm{IQR},\; Q_3 + 1.5\,\mathrm{IQR}\right],
\label{eq:tukey-fences}
\end{equation}
calculated independently within each cohort. Similar non-parametric thresholds have been used to identify parameter estimation errors in non-linear growth modelling \cite{markovic2014NewApproach}. Because this fence captures distributional extremes rather than demonstrably erroneous records, flagged plots remain in the primary analysis. A sensitivity check recalculates pooled held-out metrics after removing these observations from saved predictions, without refitting the models.

The audit evaluates whether specific plots recur within the worst 1\% and 5\% of residuals, independently across Elastic Net, XGBoost, and GNNWR. It compares their within-plot residual range, GNNWR local coefficients, edge distance to compartment boundaries, sub-compartment clustering, and topographic wind shelter against the wider sample. To test a known survey-boundary issue, the worst-fitting survey year is identified, and each flagged plot is refitted after excluding its 2008 observation. These diagnostics assess whether errors originate from target construction rather than genuine ecological factors.

%------------------------------------------------
\section{Q3: Prediction and physics guidance}
\label{sec:rq1-method}
\label{sec:q3-method}

Q3 is supplementary: it returns to whether the attribution signal in Q1 and Q2 also
improves top-height prediction, and what a physics-informed architecture does with it.
It compares a biological curve, standard tabular models and neural models. The headline
comparison uses Set~3 and the same five spatial folds. The other feature sets test whether adding
environmental inputs changes the result.

\subsection{Chapman--Richards (CR) baseline}
\begin{equation}
H(a) = y_{\max}\left(1 - e^{-ka}\right)^p,
\label{eq:cr}
\end{equation}
where $H(a)$ is top height at age $a$, $y_{\max}$ is the asymptotic height ceiling, $k$ is the growth rate parameter, and $p$ governs curve shape. All three parameters ($y_{\max}$, $k$, $p$) are fitted freely by non-linear least squares across five initial starting points, retaining the solution with the lowest sum of squared errors -- unlike Q2's plot-level fit, which fixes the equivalent rate and shape parameters at Forestry Commission values (Section~\ref{sec:q2-method}). Parameters are calibrated exclusively on the training partition of each split and fold, then frozen. This prevents validation and test observations from influencing the baseline fit or the PINN loss terms derived from it.

\subsection*{Conventional models}
Three classical baselines represent different levels of complexity: ordinary linear regression, a bagged random forest, and gradient-boosted decision trees via XGBoost \cite{chen2016XGBoost_C4_NN}. All conventional models use the same Q3 feature sets and cross-validation folds.

For Q3, XGBoost hyperparameters were tuned specifically on the top-height target. Evaluating 27 candidate configurations across five folds and two cohorts produced 270 fits. Model selection relied strictly on validation $R^2$, keeping test sets untouched until final evaluation. Appendix~\ref{app:hyperparameters} lists the full search grid and chosen values. This tuned configuration is specific to Q3; Q1 and Q2's attribution models use the fixed configuration given in Section~\ref{sec:q1-method} instead, without further retuning.

\subsection{DNN and physics-informed models}
All neural architectures share a three-layer prediction trunk with 128 units per layer and Leaky-ReLU activations. Training uses the Adam optimiser with a base learning rate of $10^{-4}$, weight decay of $10^{-5}$, gradient clipping at 1.0, batch size 256, and early stopping with a patience of 40 epochs (up to 500 epochs total). A ReduceLROnPlateau scheduler (decay factor 0.8, patience 15 epochs) lowers the learning rate when validation loss stalls, remaining active across most of the optimisation trajectory. The standard DNN minimises mean squared error:
\begin{equation}
\mathcal{L}_{\mathrm{data}}=\frac{1}{N}\sum_{i=1}^{N}
\left(\hat H_i-H_i\right)^2.
\label{eq:data-loss}
\end{equation}

The environment-conditioned DNN concatenates the environmental variables to the main input. The PINN variants use the same prediction trunk, but environmental variables enter small auxiliary networks that shape the physics and trajectory losses (Equations~\ref{eq:physics-loss}--\ref{eq:trajectory-loss}), not the trunk's prediction itself. The architectures therefore share the same core capacity, while the environmental information reaches training through a separate pathway from the DNN's.

\limit{This separation between the trunk's prediction and the auxiliary networks was not designed to isolate a growth-ceiling-specific effect; it reflects a limitation in the implemented forward pass, identified after training and discussed in Section~\ref{sec:rq1-pinn-limitation}.} One consequence is diagnostically useful regardless of intent: because terrain never reaches the trunk, there is nothing terrain-derived for a random split to leak between train and test, which separates this architectural property from the standard DNN baseline's own known split sensitivity -- early exploratory plot-level experiments showed clear leakage for the DNN baseline, motivating the spatial-block design used throughout.

The analytical derivative of the CR curve is:
\begin{equation}
H'(a)=y_{\max}pk e^{-ka}\left(1-e^{-ka}\right)^{p-1}.
\label{eq:cr-derivative}
\end{equation}
The physics loss penalises discrepancies between this theoretical growth rate and the network derivative computed via automatic differentiation:
\begin{equation}
\mathcal{L}_{\mathrm{phys}}=\frac{1}{N}\sum_{i=1}^{N}
\left(\frac{\partial\hat H_i}{\partial a_i}-H'(a_i)\right)^2.
\label{eq:physics-loss}
\end{equation}
To maintain consistent units ($\mathrm{m\,yr^{-1}}$), the network derivative is divided by the training-set age standard deviation before computing the loss.

The trajectory loss uses consecutive training surveys from the same plot:
\begin{equation}
\mathcal{L}_{\mathrm{traj}}=\frac{1}{M}\sum_{j=1}^{M}
\left(\frac{\hat H_{j,2}-\hat H_{j,1}}{\Delta a_j}-H'(\bar a_j)\right)^2,
\label{eq:trajectory-loss}
\end{equation}
where $\bar a_j$ is the midpoint age. Pairs are formed only when both rows belong to the training
partition. The total loss is
\begin{equation}
\mathcal{L}=\mathcal{L}_{\mathrm{data}}+
\lambda_{\mathrm{phys}}\mathcal{L}_{\mathrm{phys}}+
\lambda_{\mathrm{traj}}\mathcal{L}_{\mathrm{traj}}+
10^{-5}\sum_{\theta}|\theta|.
\label{eq:total-loss}
\end{equation}
The default uses $\lambda_{\mathrm{phys}}=\lambda_{\mathrm{traj}}=1$. Joint weights 0, 1 and 2
are compared in the ablation. Weight 0 is the physics-disabled control.

The first conditioned PINN changes the asymptote:
\begin{equation}
y_{\max}^{(i)}=y_{\max}+f_y(\mathbf{x}_i).
\label{eq:ymax-adjust}
\end{equation}
PINN-$k$ also changes the rate while keeping it positive:
\begin{equation}
k^{(i)}=k\exp\!\left(f_k(\mathbf{x}_i)\right).
\label{eq:k-adjust}
\end{equation}
Both auxiliary networks have one 16-unit hidden layer and start near zero. The shape parameter
$p$ remains fixed in every variant.

The trunk's prediction never evaluates Equation~\ref{eq:cr} with $y_{\max}^{(i)}$ or $k^{(i)}$;
these enter only the physics and trajectory losses (Equations~\ref{eq:physics-loss}--
\ref{eq:trajectory-loss}). Environmental information therefore reaches the trunk only indirectly,
through those loss gradients, never as a direct input to the prediction itself. \limit{This is a
limitation of the implemented forward pass rather than a chosen ablation design; see
Section~\ref{sec:rq1-pinn-limitation}.}

\subsection*{Does environmental conditioning correct the curve at row level?}
\label{sec:rq2a-method}
This is a row-level, supplementary check, not this chapter's headline attribution result: it asks whether environmental conditioning corrects a plot--survey's error at all, as a bridge into Q1 and Q2's plot-level attribution models above. It evaluates whether environmental predictors explain departures from standard yield curves using held-out predictions from Q3's own models. For each test observation $i$, the absolute error reduction relative to the fold-matched Chapman--Richards baseline is
\begin{equation}
d_i=\left|H_i-\hat H_i^{\mathrm{CR}}\right|-
\left|H_i-\hat H_i^{\mathrm{model}}\right|.
\label{eq:residual-reduction}
\end{equation}
A positive value indicates that the model improves on the baseline prediction. The analysis reports mean $d_i$, its percentage of mean Chapman--Richards error, and the percentage of observations where $d_i>0$.

Within each held-out fold, test observations are ranked by absolute baseline error and divided into quartiles (Q1 to Q4). Q1 contains observations where the Chapman--Richards curve fits best, while Q4 contains those with the largest departures. Mean error reduction and the percentage of improved plots are calculated for each quartile. This is repeated for the environment-conditioned DNN, PINN and PINN-$k$ across Sets~2--4 and both cohorts, and for XGBoost on Set~3 and both cohorts. The Set~3 DNN is evaluated across five training seeds to measure parameter initialisation spread. Other configurations use seed 42, meaning their reported spread reflects fold-to-fold variation rather than seed variance.

Two controls, both run on XGBoost with Set~3 only, isolate the source of any improvement. First, XGBoost is evaluated with and without environmental inputs. Second, the environmental columns are permuted using a single shared row permutation across every row of a held-out fold (train, validation, and test alike), and XGBoost is refitted from scratch on the permuted data. This tests whether the performance gain requires real environment-plot linkage, avoiding the confound of evaluating an already-trained model on out-of-distribution inputs. The same error-reduction calculation is then repeated on the refitted model's held-out predictions. For the DNN and XGBoost (both Set~3, both cohorts), the spatial pattern of $d_i$ is assessed with Moran's $I$ using the same weights as the residual analysis.

%------------------------------------------------
\section{Evaluation and robustness}
\label{sec:eval-robustness}

For residual $e_i=H_i-\hat H_i$, the main metrics are
\begin{align}
\mathrm{RMSE}&=\sqrt{\frac{1}{N}\sum_i e_i^2}, &
\mathrm{MAE}&=\frac{1}{N}\sum_i|e_i|, \\
R^2&=1-\frac{\sum_i e_i^2}{\sum_i(H_i-\bar H)^2}, &
\mathrm{Bias}&=\frac{1}{N}\sum_i e_i.
\end{align}
Positive bias denotes under-prediction. Each table states whether it reports a fold mean or a
pooled metric. Fold summaries are the mean $\pm$ sample SD across five folds. Pooled metrics are
calculated once from all out-of-fold predictions. These are not interchangeable.

For the four-survey spatial analyses, pooled $R^2$ 95\% confidence intervals use 1,000 cluster bootstrap resamples of whole compartments. The 2.5th and 97.5th percentiles form the interval. This preserves spatial clustering better than resampling individual rows. The smaller six-survey Q2 results are reported as point estimates; GNNWR is also repeated over three split seeds to test sign stability.

A semivariogram measures how quickly residuals become less similar as distance between plots increases; its range is the distance beyond which they are no longer spatially correlated. Each model's range is fitted separately from its own residuals, searched up to 5000m by default; this window does not widen automatically. One diagnostic only, the six-survey Q2 comparison, was manually rerun with the window widened to 10,000m -- even then several cells remained unresolved at the new cap, reported as provisional, not as a resolved range. Two models' Moran's $I$ are not necessarily tested at the same distance, and any report of a Moran's I value states the distance it was computed at for this reason.

Global Moran's $I$ is calculated from held-out residuals using distance-band weights, with the threshold set by this semivariogram range, and significance assessed via 999 permutations (ranges used are tabulated in the Appendix):

\begin{equation}
I=\frac{N}{W}
\frac{\sum_i\sum_j w_{ij}(e_i-\bar e)(e_j-\bar e)}
{\sum_i(e_i-\bar e)^2}.
\label{eq:morans-i}
\end{equation}
It is used after fitting, not for feature selection.
Three outcomes occur: the range resolves within the search window (Moran's $I$ reported directly); the semivariance stays flat throughout the window (no spatial structure detected, Moran's $I$ not computed); or the semivariance is still rising at the window's edge (Moran's $I$ computed at that capped distance, reported as provisional, not equivalent to a resolved result). The same rule is applied to Q3's bridge check, Q1's residuals and the Elastic Net, XGBoost and GNNWR residuals in Q2.

Neural training variation is assessed over five seeds for the selected Q3 models. XGBoost and the
DNN are also compared fold by fold on identical test folds. The physics-weight ablation tests
weights 0, 1 and 2. For PINN-$k$, parameter separability is assessed by the pooled correlation
between learned $y_{\max}$ and $k$; its confidence interval uses a separate 2,000-resample compartment
bootstrap (not the 1,000-resample procedure used for the $R^2$ intervals above). These checks distinguish predictive accuracy from parameter stability.

\section{Reproducibility}

Split assignments, feature lists and model settings are saved with each run. Seed 42 is used unless
a stated reseed is part of the analysis. Full feature membership, hyperparameter searches,
secondary split results and detailed outlier tables are retained in the Appendix. Methodological
limitations and their effect on interpretation are evaluated after the results, before the final
conclusions.

%=======================================================================================
%=======================================================================================
\chapter{Results and Discussion} \label{ch:results}

This chapter answers three questions, in order: Q1 attributes a plot's persistent
departure from the shared growth curve to environmental and management conditions,
and tests whether independent model families agree. Q2 tests whether letting that
relationship vary spatially explains what Q1's global models cannot. Q3, supplementary,
tests whether any of this signal also improves top-height prediction, and reports what
a physics-informed architecture did and did not achieve for this data.

%=======================================================================================
\section{Q1: Persistent curve-departure attribution}
\label{sec:results-rq2}
\label{sec:results-q1}

\subsection{What explains a plot's persistent departure from the shared curve?}
\label{sec:results-rq2b}
\label{sec:results-q1-attribution}

4survey only. Plot-level mean departure, full survey history (Table~\ref{tab:results-rq2b}).
Result: canopy cover + thinning history have the largest Elastic Net coefficients and
top XGBoost SHAP rank -- of all predictors, every feature set, both methods.

\begin{table}[!htbp]
\centering
\scriptsize
\renewcommand{\arraystretch}{1.2}
\begin{tabular}{lrlll}
\toprule
Set & n rows & EN R2 [95\% CI] & XGB R2 [95\% CI] & LMM spatial variance explained \\
\midrule
Set1 (baseline only) & 71{,}766 & 0.220 [0.180, 0.252] & 0.222 [0.183, 0.252] & 0.016$\pm$0.078 \\
Set2 & 71{,}766 & 0.359 [0.310, 0.396] & 0.416 [0.355, 0.461] & 0.204$\pm$0.066 \\
Set3 & 71{,}727 & 0.325 [0.272, 0.365] & 0.375 [0.321, 0.414] & 0.049$\pm$0.048 \\
Set4 & 71{,}330 & 0.358 [0.305, 0.398] & 0.395 [0.338, 0.438] & 0.175$\pm$0.064 \\
\bottomrule
\end{tabular}
\caption{\small{Attribution of the persistent mean curve residual (Equation~\ref{eq:mean-cr-residual}),
4survey only, spatial-block split.}}
\label{tab:results-rq2b}
\end{table}

Canopy cover coefficient: largest or near-largest EN, every set (Set4: +1.943). SHAP
rank: 1st, every XGBoost set (1.375--1.442).
Risk: canopy cover comes from the same LiDAR pipeline as the height targets --
could just be the target restated in different units.
Checked: Pearson correlation, canopy cover vs.\ all 3 of this project's own targets.
Result: 0.35--0.47 -- moderate. Near 1.0 would mean same-variable-in-disguise; this
does not.
Checked: drop canopy cover from Set4, refit. Result: R2 drops about a third (EN
0.350$\to$0.231; XGBoost 0.388$\to$0.249) -- bigger than its SHAP margin alone
(1.2--1.7x the next variable) predicts.
So what: canopy cover carries real signal that nothing else in Set4 backfills.
Checked: does \texttt{dist\_to\_road} (the variable most likely to absorb the gap)
pick up the slack? Result: no -- its own SHAP value moves the wrong way (1.003$\to$
0.869).
Resolved: rules out simple correlation-driven credit transfer as the mechanism.
Remaining: why canopy cover carries this much unique signal -- not explained here.
Read: treat canopy cover + thinning as baseline stand-structure controls, expected
to dominate the coefficients by construction. Look past them for the true
environmental-variable question.

Abiotic variables: only \texttt{slope\_degrees} stable across both methods -- LMM
+0.894/+0.861, EN +1.230/+0.993 (Set3/Set4), same sign, comparable size, both.
\texttt{topex}: LMM strong and stable (+1.316/+1.323). EN weak and sign-flipping
($-0.175$ Set3, $+0.063$ Set4, SD as large as the coefficient itself).
So what: genuine disagreement between the two methods on \texttt{topex}'s effect --
not a shared, confirmed signal.

Candidate explanation 1, multicollinearity. Checked: VIF, every unstable variable.
Result: \texttt{tas\_mean} 3.44, \texttt{soilgrids\_ph} 2.12, \texttt{topex} 2.37 --
all lower than the stable \texttt{chelsa\_bio12\_precip\_mm}'s VIF of 3.91.
Resolved: ruled out. VIF size does not track which variables are unstable.

Candidate explanation 2, climatology timing. Background: CHELSA's climate variables
are a static 1981--2010 average (a genuine 30-year climate normal, verified against
the actual download source, not a naming issue) -- survey years are 2008--2023, a
real mismatch, if the unstable variable is actually one of CHELSA's.
Checked: which named variable is actually the static-CHELSA one, and which is
unstable. Result: \texttt{chelsa\_bio12\_precip\_mm} is the static 1981--2010
variable -- and it is stable, not unstable. \texttt{tas\_mean}, the unstable one, is
a different source (HadUK-Grid), averaged over each cohort's own real survey years,
not the 1981--2010 baseline.
Resolved: wrong mechanism, wrong variable -- ruled out.
So what: static vs.\ non-static does not separate stable from unstable variables at
all. \texttt{chelsa\_bio12\_precip\_mm}: static and stable. \texttt{tas\_mean}:
non-static and unstable. \texttt{topex}, \texttt{soilgrids\_ph}: static and
unstable too. Not the deciding factor.
Remaining: \texttt{soilgrids\_ph} is a genuinely static soil property -- no
survey-year question even applies to it -- still no mechanism identified at this
point.

Candidate explanation 3, spatial confounding. LMM controls for grouping with a
compartment random intercept, EN does not. When a variable varies mostly between
compartments, EN can confuse its real effect with unmeasured spatial factors.
Checked: between-compartment variance share (ICC), all 4 named variables. Result:
\texttt{topex} 0.74, \texttt{soilgrids\_ph} 0.77, \texttt{tas\_mean} 0.89 -- but the
stable \texttt{chelsa\_bio12\_precip\_mm} highest at 0.97.
So what: variance share alone does not separate stable from unstable. High
between-compartment variance is a precondition for confounding, not sufficient to
cause it by itself.

Sharper test: LMM is also linear -- one coefficient per variable, no curve. If a
variable's raw shape were the reason EN is unstable, LMM should be unstable too. It
is not. The only structural difference is LMM's random intercept. So a raw,
population-wide wavy shape is not necessarily a mechanism separate from
confounding -- it may be what confounding looks like in a pooled bivariate plot.
Checked directly: subtract each compartment's own mean from both the variable and
the target, redo the decile/correlation check on these within-compartment
residuals -- mechanically close to what a random intercept removes.
Result, \texttt{topex}: within-compartment Spearman 0.21, \emph{stronger} than the
raw 0.09, and close to monotonic (the wave mostly disappears). Real local signal
exists; EN just never sees it directly, because it fits the raw pooled relationship
instead.
Result, \texttt{tas\_mean}: within-compartment Spearman 0.08, much \emph{weaker}
than the raw 0.27. Most of its apparent signal is a between-compartment difference,
not a real local relationship.
Result, \texttt{soilgrids\_ph}: within-compartment Spearman 0.05, close to zero
(raw was 0.27). Same pattern -- little genuine local signal survives.
Result, \texttt{chelsa\_bio12\_precip\_mm} (stable): within-compartment Spearman
$-0.11$, weaker than its raw $-0.36$ but still real, and still a clean, monotonic
shape.
So what: not two independent mechanisms (confounding, then separately shape). One
mechanism, three outcomes. \texttt{topex}: real local signal, masked by
between-compartment pooling in EN's raw fit -- LMM's random intercept recovers it,
hence stable there. \texttt{tas\_mean}, \texttt{soilgrids\_ph}: almost no real
local signal at all -- their raw correlation is mostly a between-compartment
artifact, so there is little genuine relationship for any model to stably estimate
a coefficient from. \texttt{chelsa\_bio12\_precip\_mm}: real, clean local signal
survives alongside very high between-compartment variance -- enough to anchor a
stable EN coefficient despite the confounding.
Resolved: the best-evidenced account so far, and a more mechanistically unified one
than the separate confounding/nonlinearity framing this replaces. Not a full
multivariate test -- raw bivariate/variance-share checks, not the actual joint
Set3/Set4 EN fit itself.
Saved: full script, binned tables, within-compartment decomposition, caveats --
\texttt{TEMP\_results\_attribution/TEMP\_rq2b\_topex\_tasmean\_nonlinearity\_check\_2026-08-20.tex}
(appendix candidate).

Result: residual spatial clustering real and significant, every set (distance-band
Moran's $I$, resolving at 1385--1739\,m, $p=0.001$).
Result: EN's own clustering barely moves, Set1$\to$Set4 (0.145$\to$0.103, $-29\%$
relative).
Result: XGBoost's own clustering drops much more, same sets (0.146$\to$0.054,
$-63\%$ relative).
So what: real divergence between the two attribution methods in how much spatial
structure each removes -- not the same, near-identical pattern a coarser check
would suggest.
%CHECK: verified against model and ecological literature -- supported model theory,
%cause untested.
Why (theory): XGBoost learns nonlinear thresholds and interactions automatically;
EN remains linear and additive unless nonlinear or interaction terms are added
manually \cite{zou2005Regularization_C4_NN}\footnote{XGBoost feature interaction
constraints documentation: \url{https://xgboost.readthedocs.io/en/stable/tutorials/feature_interaction_constraint.html}}.
Correction: not specifically because predictors are ``spatially smooth''; because
environmental--growth relationships may themselves be nonlinear or interactive.
Possible mechanism (untested): XGBoost better represents thresholds and
interactions among elevation, exposure, moisture, climate, and soil conditions.
Alternative explanation (untested): greater flexibility also lets XGBoost fit
spatial proxies or noise; a lower Moran's $I$ does not by itself prove better
ecological representation.

%CHECK: verified against spatial-model theory.
Remaining clustering: evidence that predictions retain spatially patterned error.
Possible causes: omitted environmental drivers; environmental data resolution;
management history; genetics; disturbance; measurement error; spatially varying
relationships \cite{gaspard2019Residual_C5_EVAL}.
Correction: ``spatial autocorrelation remains'', not ``spatially structured
variance remains''.
Implication: supports testing a spatially varying model; does not prove spatially
varying coefficients are the cause.

Result: Set3 (terrain-heavy, no soil or climate coverage) is the weakest set on
both point estimate and interval position.
Checked: category refit ablation, Set4's actual 19 columns (not the earlier,
differently-scoped exploratory notebook), single spatial-block split. Each
category removed in turn, model refit, R2 compared against the full Set4 model.
Result: \texttt{soil\_site} helps, modestly ($+0.022$ R2 when present).
\texttt{climate} hurts ($-0.056$ R2 when present -- removing it \emph{raises} R2).
\texttt{terrain} and \texttt{spatial\_position\_edge\_effects} are Set4's two
largest positive contributors, both bigger than soil.
So what: ``soil and climate together explain Set4's edge over Set3'' is not
supported. Soil contributes a small real gain; climate is actively harmful here,
not helpful.
Correction: forest site productivity genuinely depends on combined biophysical
conditions in general \cite{bontemps2014Predictive_C2_BG}, but that does not mean
every added category helps in this specific data -- climate's own instability
(established above) is a plausible reason it hurts more than it helps here.
Remaining: why terrain and edge-effects, not soil and climate, carry most of
Set4's real gain -- not the pair originally guessed, still not fully explained.
Saved: full ablation, method, caveats --
\texttt{TEMP\_results\_attribution/TEMP\_rq2b\_set4\_category\_ablation\_2026-08-20.tex}
(appendix candidate).

%CHECK: statistical interpretation.
Result: every set's R2 confidence interval clears Set1's (baseline-only) interval,
no overlap.
Non-overlapping intervals: strong evidence consistent with environmental benefit; a
paired difference interval would still be the preferable test.
Result: Set2 vs.\ Set4 (XGBoost) intervals overlap substantially ([0.355, 0.461]
vs.\ [0.338, 0.438]).
Overlapping intervals: ranking uncertain visually; not evidence the two sets are
equivalent.
Moran's $I$ reduction (EN vs.\ XGBoost, above): descriptive evidence of method
divergence; a formal statistical difference between the two has not been tested.

\textbf{For Q1}, structural and management variables explain most of
the attributable signal in a plot's persistent departure from the shared curve, and
this is not simply the height targets restated in different units: checked directly,
and ruled out. Among true environmental variables, only slope shows a reliably
consistent association across independent model families; the others destabilise for
a specific, evidenced reason. Elastic Net fits each variable's raw, pooled
relationship to the target; LMM's compartment random intercept removes
between-compartment differences first. For \texttt{topex}, that pooling masks a
real, cleaner local signal LMM recovers and EN does not, which is why LMM is stable
there and EN is not. For \texttt{tas\_mean} and \texttt{soilgrids\_ph}, little real
local signal survives the same test at all -- most of their apparent association is
a between-compartment difference, not a genuine relationship either model could
stably estimate. \texttt{chelsa\_bio12\_precip\_mm}, the one stable predictor
checked this way, keeps a real local signal on top of the highest between-compartment
variance of any variable tested -- confounding alone does not explain instability;
whether real, local signal survives it does.
At the category level, environmental contribution is uneven even among variables
that do carry some signal: soil adds a small real gain within the richest feature
set, while climate is actively counterproductive there. Elastic Net and XGBoost
diverge sharply in how much residual spatial clustering they remove (29\% vs.\ 63\%),
and a substantial share of the departure's spatial pattern is not explained by any
environmental set tested here regardless of method -- the direct motivation for
testing a spatially-varying model in Q2.

\fig{\textbf{Q1 -- coefficient/Moran's I/residual-map composite.} Three panels.
Panel A (attribution): coefficient forest plot, one row per variable (\texttt{CanopyCover},
both thinning columns, \texttt{slope\_degrees}, \texttt{topex}, the climate/soil
variables), point = mean, whisker = fold SD, LMM and EN as two colours/markers side by
side, Set3 and Set4 as two sub-columns -- the visual payoff of the slope-stable/
topex-unstable contrast, otherwise the reader has to hold four numbers in their head.
Panel B (spatial clustering responds unevenly to environmental data): Moran's I by set
(Set1--Set4), EN and XGBoost as two lines -- EN nearly flat, XGBoost dropping
substantially, the divergence between the two lines is the finding, not a shared trend.
Panel C (where the leftover residual sits): a map of Set4/4survey's model residual,
coloured by magnitude/sign -- this is what motivates Q2, earning its panel on
structural grounds, not decoration. Establishes: none of this (coefficients, the
Moran's I trend, the spatial residual pattern) exists in any table in this chapter, all
of it is currently text-only. Answers Q1 items 1--3.}

\fig{\textbf{Q1 -- two variables, equally spatially confounded by eye.} Two-panel map
of Aberfoyle, same colour treatment on both panels for direct visual comparison.
Panel 1: \texttt{topex} value per plot, continuous colourmap, compartment boundaries
drawn as thin outlines. Panel 2: \texttt{chelsa\_bio12\_precip\_mm}, identical
treatment. Both should read as visibly blocky, colour aligning with compartment
boundaries -- the visual counterpart to their ICC numbers (0.74 and 0.97). Sets up a
question the next figure answers: these two look equally confounded by eye, so why
does only one destabilise the Elastic Net coefficient? \texttt{tas\_mean} and
\texttt{soilgrids\_ph} deliberately left out -- also high-ICC, so their maps would
look similar without adding a new visual point; their actual story (real signal
nearly vanishes) belongs to the next figure instead. Establishes: the geometric
precondition for spatial confounding, visually, not just as a number. Answers Q1
item 3 (mechanism groundwork).}

\fig{\textbf{Q1 -- does each variable's signal survive removing compartment
structure?} Dumbbell/slope chart. Y-axis: signed Spearman correlation with
\texttt{mean\_cr\_residual\_4survey}, zero line marked. Four rows, one per named
variable (\texttt{topex}, \texttt{tas\_mean}, \texttt{soilgrids\_ph},
\texttt{chelsa\_bio12\_precip\_mm}). Two points per row, connected by a line:
``raw'' (population-wide) and ``within-compartment'' (compartment means removed
from both variable and target before correlating). Colour or label each row by EN
stability (stable/unstable) so the up/down direction visually maps onto the actual
outcome. Values: topex $+0.093\to+0.210$ (up -- real local signal was masked);
tas\_mean $+0.273\to+0.082$ (down toward zero); soilgrids\_ph
$+0.271\to+0.046$ (down toward zero); chelsa\_bio12\_precip\_mm
$-0.364\to-0.114$ (down, but stays clearly non-zero). No variable flips sign --
direction is stable within each row, only magnitude changes, so the up/down
reading is not an artefact of sign choice. \%SOWHAT: chelsa\_bio12\_precip\_mm's
negative sign (more rain, shorter-than-expected growth) runs against Sitka
spruce's documented rainfall preference -- untested whether precipitation itself
or something it proxies for (elevation, exposure, drainage) is doing the work; not
resolved here. Establishes: the actual mechanism the previous figure's question
sets up -- confounding is necessary but whether real local signal survives it is
what decides EN's stability. Answers Q1 item 3 (mechanism, resolved).}

%TODO: checked and changed 2026-08-20 -- topex actually has 16,839 distinct values (verified
%directly against data/processed/environmental/plot_environmental_features.parquet), not
%13--269. soilgrids_ph (13), tas_mean_4survey (149), chelsa_bio12_precip_mm (269) do fit that
%range; topex is the outlier and was wrongly folded into it. NOT a pure typo fix: the original
%framing ("distinct values" as a proxy for "coarse grid vs. plot-level") is complicated by this,
%since topex is fine-grained (more distinct values than tas_mean) but still highly confounded
%(ICC 0.74) -- distinct-value count alone does not predict spatial confounding, ICC does. This
%arguably sharpens the appendix's own point (CanopyCover's low ICC isn't just "it has more
%distinct values," since topex also has many and is still confounded) but changes what the
%figure caption should claim. Flagging for the other chat / a second pass -- not resolved
%unilaterally here, just corrected to not state a false number in the meantime.
\fig{\textbf{Q1 (Appendix) -- a low-confounding contrast.} Single map,
\texttt{CanopyCover} per plot, same colourmap/compartment-boundary treatment as the
topex/chelsa\_bio12\_precip\_mm figure above, for direct visual contrast. ICC 0.30
(vs.\ 0.74--0.97 for the four environmental variables checked) and $\sim$70,769
distinct values (vs.\ 13--269 for soilgrids\_ph/tas\_mean/chelsa\_bio12\_precip\_mm;
topex separately at 16{,}839 -- fine-grained but still highly confounded, ICC 0.74) --
a plot-level, management-driven variable rather than a coarse grid one. Should read as visibly speckled within each compartment,
not blocky -- the calibration reference the main-text map figure needs but, with
only two high-ICC examples, cannot provide on its own: what does genuinely low
spatial confounding actually look like, next to what topex and
chelsa\_bio12\_precip\_mm look like. Appendix, not main text -- supporting/
calibration material, not part of the load-bearing argument. Establishes: visual
grounding for reading the two main-text figures correctly.}

%=======================================================================================
\section{Q2: Spatial attribution}
\label{sec:results-rq3}

Q1 result: global models leave real spatial clustering unexplained; XGBoost erased
far more of it than EN as environmental data was added.
Q2 tests: does letting the environment-deviation relationship vary spatially --
not one global relationship -- capture what the global models could not
(Table~\ref{tab:results-rq3})?

\begin{table}[!htbp]
\centering
\scriptsize
\renewcommand{\arraystretch}{1.2}
\begin{tabular}{lllllll}
\toprule
Set & EN 4survey [CI] & EN 6survey & XGB 4survey [CI] & XGB 6survey & GNNWR 4survey [CI] & GNNWR 6survey \\
\midrule
Set2 & 0.286 [0.241, 0.325] & 0.057 & 0.266 [0.214, 0.308] & 0.031 & 0.320 [0.270, 0.369] & $-0.044$ / 0.014$\pm$0.053 \\
Set3 & 0.274 [0.223, 0.316] & 0.031 & 0.281 [0.235, 0.320] & 0.086 & 0.319 [0.263, 0.365] & 0.058 / 0.054$\pm$0.055 \\
Set4 & 0.240 [0.192, 0.286] & 0.056 & 0.250 [0.201, 0.295] & 0.015 & 0.294 [0.236, 0.347] & $-0.075$ / 0.002$\pm$0.068 \\
\bottomrule
\end{tabular}
\caption{\small{Attribution of the plot-specific asymptote deviation
(Equation~\ref{eq:local-ymax-diff}), spatial-block split. 6survey GNNWR column shows
seed-42 point estimate / three-seed mean$\pm$SD.}}
\label{tab:results-rq3}
\end{table}

Result: GNNWR's point estimate beats both global models, every 4survey set --
consistent in direction and magnitude, not a one-set fluke.
Checked: EN/XGBoost's own bootstrap confidence intervals vs.\ GNNWR's. Result:
overlap, every set (Set4: GNNWR [0.236, 0.347] vs.\ EN [0.192, 0.286]).
So what: the point-estimate edge is not confirmed at the strictest statistical bar.

Checked: independent diagnostic, residual Moran's $I$, GNNWR vs.\ the best global
model, 4survey, all 3 sets. Result: GNNWR reduces it in every set, by a larger
relative margin than the R2 gap alone suggests (Set2 $-32.0\%$, Set3 $-24.7\%$,
Set4 $-23.7\%$; all ranges resolving 808--1252\,m, $p=0.001$).
So what: corroborates the same direction as the R2 gap without resolving the CI
overlap -- evidence GNNWR's advantage reflects real spatial structure, not just a
better mean fit.

Checked: same diagnostic, 6survey. Result: uninformative, not contradictory --
almost no cell resolves a semivariogram range at all, so this cannot corroborate
or contradict GNNWR's near-zero, sign-unstable 6survey R2; that instability
stands on its own evidence, not doubly confirmed.

Candidate explanation: reference-point sparsity -- too few nearby training plots
for GNNWR's spatial-weighting network to learn a stable local pattern from, a
recognised limitation of classical GWR bandwidth selection too
\cite{brunsdon1996Geographically_C6_GWR}. 6survey has far fewer compartments than
4survey (48 vs.\ 296).
Checked: reference-density ablation, GNNWR refit on 4survey, current Set3/Set4
scope, reference-set sizes from 7,500 (approximating 6survey's own scale) up to
full, single spatial-block split.
Result: 4survey's R2 stayed positive and broadly similar at every size tested,
including the smallest (Set3 0.27--0.34, Set4 0.12--0.19) -- never approaching
6survey's near-zero, sign-flipping range.
Resolved: reference sparsity alone does not reproduce 6survey's collapse -- not
sufficient by itself.
Remaining: two things limit how far this rules the account out. R2 did not move
monotonically with reference size (Set4's highest value came at the smallest size
tested, unexplained, possibly single-split noise). Having fewer points alone does
not explain the collapse: the ablation samples fewer points but still
proportionally from every compartment, so shrinking it this way does not capture
6survey's tighter geographic spread, which comes from having far fewer
compartments to begin with, not just fewer points spread over the same area.

Result: the same $\sim$10 plots recur as the worst residual, nearly every feature
set and cohort combination -- rules out a poor choice of features as the
explanation.
Checked, scaled to the full population below: does the recurrence hold across
model families, not just feature sets? Result: yes -- flagged plots land in each
model's own worst-1\%/worst-5\% $|$residual$|$ bucket at roughly 36x the chance
rate, independently in EN (35.7\%/59.0\%), XGBoost (38.0\%/61.3\%), and GNNWR
(37.6\%/57.5\%).
So what: three structurally different methods agree on which plots are hard to
fit -- not an artefact of one method's own quirks.

Checked: disturbance flags (clearfell, measurement-anomaly, trajectory-instability),
the 10 plots. Result: 0 of 10 trigger any flag.
Result, trajectory shape: half smooth and internally consistent -- the large
deviation comes entirely from sitting far from the yield-class benchmark curve,
consistent with a yield-class lookup mismatch, not an observed-height problem.
Half show real, sub-threshold curve-fit instability across their own survey years.
Checked: trace 5 of the 10 plots' implausible fitted asymptotes (up to 116\,m,
against normal raw heights of 38.9--46.8\,m) back to the fitting procedure itself.
Result: confirms the mechanism directly -- the closed-form fit holds the
growth-rate parameter fixed at the plot's yield-class value, solves only for the
asymptote; a plot whose real trajectory is flatter or slower than its yield class
assumes gets forced into an inflated ceiling to reconcile the mismatch.
Resolved: a property of how the target is constructed, not a data-quality gap --
no height- or asymptote-based filter in the cleaning pipeline would have caught
it, since the raw heights are unremarkable.

Checked: GNNWR's own local coefficients, flagged vs.\ control plots. Result:
flagged plots carry $\sim$2.7x higher prediction error (21.12\,m vs.\ 7.70\,m) but
statistically indistinguishable local coefficients.
\sr{TODO -- overclaim check: this and every other local-coefficient claim in Q2
(the CanopyCover/topex ranking below, the "For Q2" synthesis, Panel C's
CanopyCover coefficient map) treat GNNWR's local coefficient as a clean read of
genuine spatial non-stationarity. GNNWR is still locally LINEAR -- one straight-line
coefficient per predictor per location, same functional-form assumption as EN, not
XGBoost. Q1's own finding that XGBoost erases more residual clustering than EN,
and Q3's finding that 27\% of XGBoost's SHAP weight sits in interactions, are
direct evidence real non-linear/interaction structure exists in this data. GNNWR
can't tell "this coefficient is unstable because the true relationship varies by
place" apart from "this coefficient is unstable because forcing a straight line
onto a non-linear relationship fits differently well in different local samples"
-- both look identical from the outside. The confound-isolation argument for
GNNWR (holding functional form fixed to isolate spatial variation as one clean
factor) defends the R2-comparison headline result; it does NOT defend claims read directly
off individual local coefficients, including this one. Revisit once Methodology's
Q2 section is finalised -- soften "finds nothing spatially unusual" and the
downstream echoes of it, or add an explicit limitation here.}
So what: a model built specifically to detect local spatial variation finds
nothing spatially unusual about these plots -- direct evidence the deviation is a
target-construction artefact, not a real environmental or spatial signal.

Result: not a handful of edge cases. A Tukey-fence check on the fitted asymptote
flags 0.47\% of 4survey (266 plots), 5.32\% of 6survey (709 plots).
Result: flagged plots also cluster tighter within their own compartments than
the compartment centroid does -- more consistent with a localised event than a
compartment-wide labelling error.
Checked: does dropping flagged plots change pooled R2? Result: barely, 4survey
($-0.002$); worse on 6survey (EN $-0.027$, XGBoost $-0.037$).
So what: excluding them would remove signal the models rely on, not clean the
data.
Checked: disturbance flags, all 266. Result: 0 of 266 clearfell-like or
measurement-inconsistent.
So what: rules out abrupt, catastrophic disturbance -- but these flags only catch
sudden drops, not chronic effects, so this does not by itself confirm every
flagged plot is a pure curve-fitting artifact.

Checked: does the fixed-k mechanism actually hold for every flagged plot
individually, not just the 5 traced examples? Compared each flagged plot's own
observed growth (first to last survey) against its yield-class benchmark curve's
predicted growth over the same span -- the direct signature of the traced
mechanism (a plot growing slower than its assigned rate gets its asymptote
inflated to compensate).
Result: 71.6\% of a recomputed flagged set (179/250, close to but not identical to
the chapter's own 266 -- a minor methodology difference, not re-derived exactly
here) fit the signature. 28.4\% (71/250) do not -- their observed growth matches
or exceeds the benchmark, so the traced mechanism does not explain their
implausible asymptote.
So what: the fixed-k account holds for a majority, not all, of the flagged plots.
Checked: environmental profile, the two subsets against each other. Result:
sharply different terrain exposure -- the artifact-fitting subset sits in exposed
terrain (median \texttt{topex} $-5.48$), the non-fitting subset in sheltered
terrain (median $+2.58$, $p=0.0002$). Thinning history differs more weakly
(non-fitting subset thinned more recently and more often, $p=0.051$).
So what: two different stories, not one. Exposed plots whose growth falls short
of their assigned yield class fit the mechanical fixed-k account. Sheltered plots
that still end up flagged do not -- a real, unusually fast trajectory in good
conditions is at least as plausible an account there as any artifact, though not
confirmed further here.
Checked: chronic wind exposure (not caught by the abrupt-drop disturbance flags)
within the artifact-fitting subset specifically. Result: a real, modest
additional association (1.40x the growth-deficit baseline rate, $p=0.0007$), not
the whole story. Full test saved:
\texttt{TEMP\_results\_attribution/TEMP\_rq3\_flagged\_plot\_mechanism\_split\_2026-08-20.tex}
(appendix candidate).

\limit{Two refinements to the disturbance flags follow directly from this
(exposure-aware reclassification; a new chronic-underperformance flag) --
implemented and tested (\texttt{models/growth\_curve\_attribution/
disturbance\_checks\_refined.py}) but not wired into the curve-fitting exclusion
pipeline, since doing so correctly would require refitting every downstream
model to keep this chapter's own numbers consistent, not feasible within the
submission timeline. Full rationale, numbers and the worst-fitting-survey-year/
2008-boundary-artifact check: Appendix~\ref{chap:appendix} and
\texttt{TEMP\_results\_attribution/TEMP\_disturbance\_flag\_refinement\_2026-08-20.tex}.}

Result: which variable ranks first depends on cohort -- canopy cover ranks first
on 4survey by three converging methods (EN coefficient, XGBoost gain importance,
SHAP), but XGBoost disagrees with EN/GNNWR on 6survey.
\sr{TODO -- overclaim check: GNNWR's side of this "converging methods" agreement
carries the same locally-linear-vs-non-linear caveat flagged at the flagged/control
local-coefficient claim above -- don't read GNNWR's agreement/disagreement with
EN and XGBoost here as fully independent corroboration without that caveat stated
somewhere in this section.}
So what: 6survey's own XGBoost R2 is near zero or negative there
(Table~\ref{tab:results-rq3}) -- an independent, non-circular reason to trust its
6survey attribution less regardless; not the primary cohort for Q2's spatial
questions given GNNWR's own established unreliability there.

Checked: predictors grouped into terrain/wind/soil/climate/spatial-position
categories (4survey, Set4). Result: stand structure dominates (roughly 3x the
next-largest category); among true environmental categories, terrain contributes
most, soil is indistinguishable from noise -- consistent with the
individual-variable pattern above. Category-level Moran's $I$ diagnostics
(including a band-width artifact found and corrected in the terrain category's
own check) are reported in Appendix~\ref{chap:appendix}, not repeated here.

\textbf{For Q2}, the evidence supports the hypothesis that the environment-deviation
relationship varies spatially, on the cohort with enough spatial reference density
to test it: reference sparsity alone was tested directly and does not explain
6survey's collapse, so the 4survey/6survey contrast is not simply an artefact of
sample size. The accuracy edge alone does not clear the strictest statistical bar,
and needs the independent Moran's $I$ result to be read as corroborated rather than
confirmed. Separately, and now agreed on by three structurally different
attribution methods rather than one, most (though checked directly and not all)
of the largest fitted deviations trace mechanically to the target's own
construction under a fixed growth-rate assumption -- GNNWR's own local
coefficients find nothing spatially unusual about the affected plots despite
their much larger error, direct evidence against a real environmental or spatial
cause for that majority.
\sr{TODO -- overclaim check: same caveat as the flagged/control local-coefficient
claim earlier in this section (GNNWR is locally linear, can't separate genuine
spatial non-stationarity from a forced-linear misfit) -- this synthesis sentence
restates that claim as settled evidence; soften or caveat once Methodology's Q2
section states the limitation properly.} The minority this mechanism does not explain sits in
markedly more sheltered terrain than the majority does, a real, unresolved
difference; and even within the majority, chronic wind exposure -- which the
disturbance flags used here structurally cannot detect, since they require an
abrupt drop -- shows a modest but genuine additional association. This should
temper how far any single extreme \(\Delta y_{\max}^{\text{plot}}\) value is
read as a real biological signal, in either direction: some are artefacts, some
plausibly are not, and this chapter does not yet know which is which for any
individual plot.

\fig{\textbf{Q2, Figure 1 -- GNNWR edge/spatial-structure/local-coefficient
composite.} Three panels. Panel A: pooled R2 by set and model (EN/XGBoost/GNNWR),
point+CI, 4survey only. Panel B: residual Moran's I by set/model, 4survey only (mark
6survey's cells as unresolved/not-significant rather than plotting a misleading point
value) -- the independent corroboration for Panel A's edge, and the reason 6survey is
absent rather than shown as a contradictory result. Panel C: a spatial map of GNNWR's
own per-plot local
coefficient for \texttt{CanopyCover} (its top-ranked variable), coloured by
coefficient value -- the only way to show what "spatially-varying coefficient" means
concretely, no prose substitutes for seeing it vary across Aberfoyle. Establishes:
Panel B is a diagnostic that exists nowhere else visually; Panel C is the one thing in
the whole chapter that only GNNWR can produce. Answers Q2 items 2--3.
\sr{TODO -- overclaim check: same caveat as the flagged/control local-coefficient
claim above. Panel C's caption/discussion should not present the mapped coefficient
as a clean read of a genuinely spatially-varying CanopyCover effect without at least
a pointer to the locally-linear limitation -- a viewer could otherwise read local
variation in this map as settled biological signal rather than a value that a
forced-linear local fit could also produce from non-linear structure.}}

\fig{\textbf{Q2, Figure 2 -- deviation map/archetype map/trajectory composite,
including the 266-plot pattern.} Three panels. Panel A: a map of
\texttt{local\_y\_max\_difference} across 4survey, coloured by deviation
magnitude/sign, with the 266 Tukey-fence-flagged plots marked with a distinct
outline/marker -- the direct visual for "the pattern scales to hundreds of plots,
clustered by compartment" (item 6), and shows at a glance whether it clusters
spatially. Panel B: the compartment archetype map (the six classes from item 9), same
geography, so a reader can cross-reference Panel A's raw deviation against Panel B's
classification directly. Panel C: small-multiples trajectories, one per archetype,
observed height vs.\ age (grey), the fitted local curve, and the yldc benchmark curve
overlaid -- the only way to make "smooth-but-offset" vs.\ "moderately unstable"
concrete rather than a label. Establishes: this figure is what lets the outlier-plot
prose (items 4--9, currently the densest part of the chapter) shorten, rather than the
other way round. Answers Q2 items 4--6, 9.}




%=======================================================================================
\section{Q3: Prediction and physics guidance}
\label{sec:results-q3}

Q1 and Q2 established that real, if partial and in places contested, environmental
attribution signal exists in a plot's departure from the shared growth curve, and that some
of it varies spatially. This supplementary section returns to the question the dissertation
opened with: does any of that signal translate into a predictive advantage for top height
itself, and does a physics-informed architecture make better use of it than a plain one?

Result: tuned XGBoost gives the most accurate top-height predictions, every model
tested, both cohorts, spatial-block evaluation. CR/physics guidance does not close
this gap at any tested loss weight ($\lambda$) -- though CR still helps parameter
identifiability (separately estimating $y_{\max}$ and growth rate $k$).

\subsection{Top-height prediction: results}
\begin{table}[!htbp]
\centering
\scriptsize
\setlength{\tabcolsep}{4pt}
\renewcommand{\arraystretch}{1.2}
\begin{tabular}{lrrrrrr}
\toprule
& \multicolumn{3}{c}{Four-survey} & \multicolumn{3}{c}{Six-survey} \\
\cmidrule(lr){2-4} \cmidrule(lr){5-7}
Model & RMSE${\downarrow}$ & MAE${\downarrow}$ & $R^2{\uparrow}$ & RMSE${\downarrow}$ & MAE${\downarrow}$ & $R^2{\uparrow}$ \\
\midrule
Linear & $5.168{\pm}0.361$ & $4.056{\pm}0.343$ & $0.580{\pm}0.026$ & $4.877{\pm}0.643$ & $3.619{\pm}0.426$ & $0.509{\pm}0.067$ \\
RF & $4.794{\pm}0.187$ & $3.611{\pm}0.142$ & $0.638{\pm}0.011$ & $4.062{\pm}0.249$ & $3.010{\pm}0.185$ & $0.656{\pm}0.046$ \\
XGBoost (tuned) & $\mathbf{4.548{\pm}0.137}$ & $\mathbf{3.431{\pm}0.123}$ & $\mathbf{0.674{\pm}0.009}$ & $\mathbf{3.656{\pm}0.232}$ & $\mathbf{2.669{\pm}0.183}$ & $\mathbf{0.722{\pm}0.031}$ \\
DNN & $4.681{\pm}0.230$ & $3.533{\pm}0.201$ & $0.655{\pm}0.016$ & $3.903{\pm}0.316$ & $2.867{\pm}0.275$ & $0.684{\pm}0.031$ \\
PINN$^*$ & $5.206{\pm}0.399$ & $4.063{\pm}0.407$ & $0.573{\pm}0.036$ & $3.886{\pm}0.290$ & $2.819{\pm}0.196$ & $0.686{\pm}0.038$ \\
PINN-$k^*$ & $5.197{\pm}0.402$ & $4.046{\pm}0.421$ & $0.575{\pm}0.037$ & $3.895{\pm}0.282$ & $2.816{\pm}0.196$ & $0.684{\pm}0.038$ \\
\bottomrule
\end{tabular}
\caption{Model comparison for top-height prediction, Set3, spatial\_block\_kfold,
per-fold mean $\pm$ SD across five folds. 4survey n=232{,}292 row-years (58{,}073
plots); 6survey n=82{,}614 row-years (13{,}769 plots). Table shows pooled mean $\pm$ SD across all folds, and margin reflects the full split rather than single folds.
$^*$PINN/PINN-$k$: the implemented forward pass does not carry environmental
information to the prediction itself, only to auxiliary loss terms -- see
``PINN: what it did and did not test'' below. Results are for the architecture
as implemented, not as intended.}
\label{tab:results-rq1}
\end{table}

\subsection{Discussion}
This section compares XGBoost against the DNN specifically -- the one neural
model implemented as intended, with environmental information reaching its
prediction directly. PINN/PINN-$k$ are addressed separately below, since their
own comparison against XGBoost is confounded by the implementation limitation.

Result: XGBoost beats the DNN, both cohorts.
Checked: fold-by-fold, not just pooled. Result: XGBoost wins 9 of 10 folds (all 5
on 6survey; 4 of 5 on 4survey, losing one by $-0.007$ $R^2$).
Checked, four candidate explanations for the margin:

\begin{itemize}
\item \textbf{Feature interactions (insufficient explanation):} XGBoost allocates 27\% of its SHAP prediction weight to interactions (primarily Age crossed with terrain and stand structure). However, DNN already captures 80\% of the nonlinear gain above the linear baseline ($R^2$ 0.580 $\to$ 0.655 out of a possible 0.674; Table~\ref{tab:results-rq1}), and Random Forest underperforms the DNN ($R^2 = 0.638$). Interaction capacity alone therefore does not explain the margin.
\item \textbf{Compartment sample size (ruled out):} Subsampling the 4-survey cohort down to 48 compartments (matching the 6-survey count) did not widen XGBoost's lead; the $R^2$ difference fluctuated between $-0.085$ and $+0.029$ across random subsets, compared to $+0.024$ on full data.
\item \textbf{Architecture-size under-tuning (ruled out):} Sweeps across five network sizes showed no consistent gains for DNN.
\item \textbf{Terrain extrapolation in held-out compartments (ruled out):} On the 4.1\% of 4-survey test plots with out-of-range terrain values, XGBoost degraded more relative to its in-range error than DNN did (error ratio 1.312 vs.\ 1.243; correlation with out-of-range distance $\rho=0.074$ vs.\ $0.059$).
\item \textbf{Training hyperparameter under-tuning (ruled out):} A 48-configuration sweep over dropout and learning rate yielded no meaningful gains for neural models under 5-fold cross-validation (within $\pm0.008$ $R^2$ of default baselines). Similarly, XGBoost's 27-configuration tuning shifted $R^2$ by only 0.002--0.004, confirming that neither model family had substantial untapped hyperparameter headroom.
\end{itemize}

Remaining: XGBoost, as a boosted-tree model, plausibly suits this tabular task
better than this DNN architecture -- consistent with, not proven by, literature
on boosted trees' strength on heterogeneous tabular data
\cite{shwartzziv2021Tabular_C4_NN}.

\fig{\textbf{Q3 -- XGBoost vs.\ DNN, fold-by-fold and the compartment-subsample
null result.} Two panels, both graphs -- nothing here is spatial. Panel A: a
slope/dumbbell chart, one line per fold (10 total, 5 folds $\times$ 2 cohorts),
connecting each fold's XGBoost $R^2$ to its DNN $R^2$ -- the fold-by-fold
consistency (9 of 10 tilt the same direction) is the actual evidence behind
``XGBoost wins,'' not just the pooled table number, and a bar chart of pooled
means would hide it entirely. Panel B: a strip/violin plot of the
XGBoost-minus-DNN $R^2$ gap across the random 48-compartment subsamples
($-0.085$ to $+0.029$), with the full-data gap ($+0.024$) marked as a
reference line -- shows directly that the compartment-count effect sits inside
ordinary resampling noise, the strongest single piece of ruled-out evidence,
shown rather than just claimed. Establishes: the two things a reader most
needs to see to trust ``XGBoost wins, and it isn't just compartment count,''
neither of which exists in Table~\ref{tab:results-rq1} or the ruled-out list
alone.}

\paragraph*{Does environmental conditioning correct the curve at row level?}
Checked: does environmental conditioning correct a row's error against the
fold-matched CR baseline (Equation~\ref{eq:residual-reduction})?
Result: yes, unevenly -- concentrates on the worst-fitting quartile (up to
5.08\,m reduction, 4survey); best-fitting quartile gets slightly worse.
Checked: permutation control, environmental columns shuffled across rows.
Result: worst-quartile correction shrinks 18--27\% but does not disappear; the
best-quartile cost appears at a similar size even with scrambled values.
So what: the worst-quartile correction needs real environmental signal, at
least partly; the best-quartile cost is a general flexibility cost, not
specific to environmental conditioning.

\paragraph*{Answer to Q3.}
For Q3: tuned XGBoost gives the strongest predictions in both cohorts; the DNN
stays competitive; CR guidance does not improve accuracy for the architecture
actually tested.

\paragraph*{PINN: what it did and did not test.}
\label{sec:rq1-pinn-limitation}
Result: PINN and PINN$_k$ underperform XGBoost and roughly trail the DNN, at
every tested physics-loss weight -- raising the weight only made predictions
worse (setting it to zero gave a real 0.05 $R^2$ gain on 4survey).
Why: a genuine implementation limitation, not evidence against physics-informed
modelling generally. The prediction trunk's forward pass never carries the
per-plot adjusted asymptote or growth rate through to the output -- only the
physics and trajectory loss terms see them -- identified after these results
were produced, not designed as a controlled ablation.
Checked, two independent diagnostics supporting this reading rather than just
asserting it: a random (non-spatial) split sharply inflates the DNN's score
(4survey: $0.831$ vs.\ $0.634$, since random splits leak spatial information)
but leaves PINN's score essentially unchanged ($0.590$ vs.\ $0.584$) -- exactly
as expected if terrain genuinely never reaches PINN's trunk to leak. Dropping
the physics loss entirely closes most, though not all, of the DNN--PINN gap on
4survey ($0.080$ down to $0.026$); the remaining $0.026$ is unresolved.
Scope: confined to how the PINN/PINN$_k$ trunk combines its sub-network
outputs -- does not affect the CR baseline, Linear, RF, XGBoost, DNN, or the
Q1/Q2 attribution models, none of which route information through this
pathway. A fix was scoped (adding the CR curve evaluated at the per-plot
parameters as an explicit prediction term) but not validated within this
dissertation's timeframe; full architectural detail and the scoped fix are in
Appendix~\ref{chap:appendix}.
\limit{The results above describe the architecture actually evaluated, not the
originally intended one -- a corrected implementation could behave
differently, in either direction, untested here.}




%=======================================================================================
%=======================================================================================
\chapter{Conclusion and Future Work} \label{ch:conc}

%=======================================================================================
%=======================================================================================
\bibliographystyle{plain}
\bibliography{mybibfile}
%=======================================================================================

\appendix
\chapter{Appendix}  \label{chap:appendix}
% \TODO: create appendix secions and labels for each section.

% ===== Q2 material cut from main text 2026-08-20, prioritisation pass -- sort into a
% ===== real appendix section later. Verbatim from the main chapter, not rewritten.
%
% -- 10-plot boundary distance (retracted at population scale, cut from main text) --
% Checked: boundary distance, these 10 plots vs. population (median 9.4m vs. 28.0m --
% ~3x closer). Checked: retested at the full 266-plot population scale. Result: holds in
% the same direction but more weakly (~2x, not 3x); one measure, distance to the forest
% perimeter, reverses sign (266-plot median further from the forest edge than typical,
% not closer). Resolved: the 10-plot estimate likely over-represents a genuinely
% edge-heavy subset -- the population-scale pattern is the more reliable one, weaker but
% still present.
%
% -- 2008 survey-boundary artifact (cut from main text) --
% Checked: worst-fitting survey year, flagged vs. population. Result: 75.6% of flagged
% plots' single worst year is 2008, vs. 47.8% population-wide.
% Why: an internally-identified compartment-boundary redraw between the 2006 and 2008
% surveys -- 71,483 new plot IDs appeared, and non-consecutive plot comparisons showed
% >20m apparent height drops in over 90% of cases, almost certainly re-delineation
% rather than real forest loss (this project's own earlier data audit, not an external
% record).
% So what: a real but minor contributor -- most of the flagged-plot pattern is not
% explained by this artefact alone.
%
% -- Terrain-removal Moran's I band-width artifact, found and corrected (cut from main
% -- text, referenced from the Methodology's Evaluation-and-robustness section) --
% Result: removing terrain (the largest true environmental category) appears at first
% to LOWER residual Moran's I (0.153 -> 0.042) rather than raise it -- the opposite of
% what removing a real spatially-structured category should do.
% Checked: why. Each category's Moran's I is measured at its own semivariogram-resolved
% band -- the full model resolves a clean 1197m band, but the terrain-removed refit's
% semivariogram never levels off inside the 5,000m window, so its Moran's I is computed
% at that 5,000m cap instead, a much wider neighbourhood that pools in far-apart,
% weakly-correlated pairs.
% Checked: recompute all 7 refits at one shared, fixed 1197m band. Result: removing
% terrain now RAISES Moran's I (0.149 -> 0.167), as does removing edge-effects
% (0.149 -> 0.174) and wind (0.149 -> 0.165); climate and stand structure barely move
% (0.149 -> 0.145 and 0.147).
% Resolved: band-width artifact, not a real effect. Terrain and edge-effects do explain
% away real spatial structure, expected direction, once compared fairly. Function
% itself confirmed correct (2026-08-19); the fix was the comparison practice, not the
% code.
%
% -- Disturbance-flag refinement, full version (shortened to a \limit{} pointer in main
% -- text) --
% Two refinements to the disturbance flags: (1) measurement_inconsistent assumes
% canopy/volume stability rules out real damage, which does not hold for wind-driven top
% breakage specifically -- an exposure-aware reclassification (high-exposure matches
% relabelled possible_wind_damage rather than excluded) would affect 22 plots total
% (18/67 4survey, 4/24 6survey intervals); (2) a new chronic-underperformance candidate
% flag (cumulative growth >=6.17m below the yield-class benchmark, no abrupt drop
% required, calibrated on 4survey and applied unchanged to 6survey) flags 3.57% of
% 4survey (8.70% of 6survey at the same threshold -- deliberately not re-calibrated,
% since 4survey is this project's primary investigation line throughout). Both are
% implemented and tested (models/growth_curve_attribution/disturbance_checks_refined.py)
% but deliberately not wired into the curve-fitting exclusion pipeline:
% measurement_inconsistent/clearfell_like currently exclude plots from the fit that
% produces both Q1's and Q2's targets, so changing that list, even for 22 plots, would
% require refitting every downstream model (Q1's EN/XGBoost/LMM, Q2's EN/XGBoost/GNNWR,
% both cohorts, every set) to keep this chapter's own reported numbers consistent -- not
% feasible within the submission timeline, and GNNWR alone is cluster-bound and
% hours-scale. The refined logic's effect on plot count is known and small (0.03% of the
% combined population); its effect on the actual reported R2/coefficient numbers is not
% known, since confirming that is exactly the rerun this decision avoids. Full
% rationale: TEMP_results_attribution/TEMP_disturbance_flag_refinement_2026-08-20.tex.
% =====

% ===== Q3 material cut from main text 2026-08-20, PINN consolidated to one paragraph --
% ===== sort into a real appendix section later. Verbatim from the main chapter.
%
% -- PINN/PINN_k intended architecture, full description (already covered in Methodology
% -- Section rq1-pinn-limitation cross-ref; redundant in Results after consolidation) --
% The PINN and PINN_k architectures were intended to let terrain and wind shape each
% plot's asymptote and growth rate (Equations eq:ymax-adjust--eq:k-adjust), with these
% adjusted parameters entering the prediction itself rather than only the physics and
% trajectory losses. In the implementation evaluated in this chapter, the prediction
% trunk's forward pass does not carry y_max^(i) or k^(i) through to the output; only the
% physics and trajectory loss terms see them (Equations eq:physics-loss--
% eq:trajectory-loss).
%
% -- Physics-weight ablation, full version (compressed into the consolidated PINN
% -- paragraph in main text) --
% Raising the physics loss weight (lambda in {0,1,2}) made predictions worse every time,
% across both models and cohorts. Setting lambda=0 raised R2 by 0.05 on 4survey (a real
% gain) and by +0.010 on 6survey (within noise). The physics loss directly hurts accuracy
% by fighting the data fit on shared weights. Dropping the physics loss cuts the 4survey
% DNN-PINN gap from 0.080 down to 0.026 (DNN 0.655 vs. unconstrained PINN_k 0.629), and
% it closes the gap on 6survey (0.719 vs. 0.684). The remaining 0.026 gap's cause is
% unresolved -- plausibly PINN_k's extra sub-network parameters, plausibly ordinary
% run-to-run variation between two structurally different networks; not settled further
% here.
%
% -- Scoped fix, full description --
% A fix was scoped -- adding the CR curve evaluated at the per-plot parameters as an
% explicit term in the prediction, with the trunk contributing a residual on top of it --
% but was not validated within the timeframe of this dissertation.
% =====

%=======================================================================================
\end{document}
%=======================================================================================

